% !TeX root = ../../Book4.tex
% 纲要标题：## 第12章　Hamilton 环面作用、凸性与 Delzant 分类
% 纲要位置：几何学计划纲要.md，第 3915 行。
% 正文状态：初稿；尚未正式审核。

% 纲要内容开始
% 
% > 本章目标与证明范围：在独立的环面作用模块中完成凸性及光滑紧辛环面流形的 Delzant 分类。像为多面体、由多面体构造、由像决定同构是三个不同证明任务。 凸性主版本取紧连通辛流形上的 Hamilton 环面作用；Delzant 分类另要求作用有效且环面维数为流形实维数的一半。
% 
% 纲要内容结束

\chapter{Hamilton 环面作用、凸性与 Delzant 分类}
\label{g4:ch:12}
\BookChapterNumber{g4:ch:12}

\begin{chapterabstract}
Hamilton 环面作用的局部权决定矩映射在轨道附近的形状，紧性又把这些局部信息组织成一个凸多面体。本章先建立等变局部标准形和临界子流形工具，证明矩纤维连通及固定点像的凸包公式。对于有效的半维环面作用，顶点权进一步构成整数基；由此得到 Delzant 多面体，并用辛约化构造相应流形，再通过边界模型和局部同构的粘合完成唯一性证明。法扇、单项式坐标与整多面体将这一分类接到代数环面簇和极化。最后的选读节把固定点权用于等变积分，并以约化类的变化计算辛体积。
\end{chapterabstract}

\begin{learninggoals}
\begin{itemize}
	\item 从完整稳定子及其切片表示建立轨道局部标准形，保留有限稳定子的作用信息。
	\item 证明矩映射分量的临界子流形非退化性、正负指标的偶性以及非空水平集的连通性。
	\item 将局部凸锥和局部连通纤维用于整体论证，证明矩像是有限个固定分支像的凸包。
	\item 核验 Delzant 条件如何保证约化自由、紧且光滑，并从模长方程恢复给定多面体。
	\item 比较边界轨道模型，消去局部同构的粘合误差，明确平移与环面重标所对应的分类自由。
	\item 由法扇写出复坐标和不变除子，说明整多面体、极化与辛形式整类的对应。
	\item 在选读部分统一角周期、Euler 权及积分归一化，计算球面和射影空间的局部化与约化体积。
\end{itemize}
\end{learninggoals}

把球面绕竖轴旋转，高度函数的每个普通水平都是一条圆轨道，两个端点则对应南北两极。高度区间记录了轨道怎样在两端收缩。对两个球面的乘积，两个高度组成一个矩映射，像成为矩形；矩形内部、边和顶点分别对应二维环面、圆和固定点。这里的多面体既记录轨道类型，也记录辛面积。

\ProseParagraphBreak

第十一章从可积系统的交换流出发，在紧正则纤维附近构造环面和作用角坐标。本章从一个已经定义在整个流形上的 Hamilton 环面作用出发，研究矩映射的全部纤维，包括稳定子增大的边界轨道。首先只要求流形紧连通，不要求作用有效或半维；在这个范围内证明凸性和纤维连通。随后加入有效性与半维条件，完成 Delzant 分类\footnote{Thomas Delzant（1961年生），法国数学家，研究辛几何中的群作用与几何群论。}：固定环面、整数格及矩映射后，一个光滑紧辛环面流形由其矩多面体确定。

\ProseParagraphBreak

证明分为三个相接的任务。局部标准形和紧性说明矩像为何是多面体；逐坐标圆作用的约化说明怎样从多面体造出流形；边界处的光滑相位及凸底空间上的粘合说明为何不会得到两个不同的模型。分类完成后，再引入法扇和复化作用，解释相同数据如何给出代数环面簇。\secref{g4:sec:1206}保持选读，它使用第二卷附录 C 的等变上同调，不参与前面分类定理的证明。

\ProseParagraphBreak

全章采用 \(T=\mathfrak t/(2\pi\Lambda)\) 和正指数基本场，矩映射满足 \(\mathrm{d}\mu^\xi=\iota_{\xi_M}\omega\)。因此正旋转权 \(\alpha\) 在标准复直线上的矩二次项是 \(-|z|^2\alpha/2\)，多面体的面则统一用原始外法向量表示。主要先修是第八章的相对 Moser 方法、第十章的正则约化和第十一章的作用角坐标。第二卷的管状邻域与普通 Morse 理论将在使用处恢复；\footnote{Harold Calvin Marston Morse（1892--1977），美国数学家，建立以临界点研究流形拓扑的方法，并研究大范围变分法。}临界子流形和负法丛所需的推广在本章证明。

% !TeX root = ../../Book4.tex
% 纲要标题：### 12.1 矩映射的局部结构与 Morse–Bott 工具
% 纲要位置：几何学计划纲要.md，第 3919 行。
% 正文状态：初稿；尚未正式审核。

% 纲要内容开始
% 
% 就地证明凸性论证需要的局部标准形和次水平集变化，引入第二卷普通 Morse 工具时明确新增的临界子流形与负法丛。
% 
% * 固定点附近的权模型与相容复法丛
% * 正则轨道与局部辛标准形
% * 临界子流形、法向 Hessian 与 Morse–Bott 引理
% * 负法丛盘丛与跨临界值的次水平集变化
% * 一般圆方向矩映射分量的 Morse–Bott 性
% * 指标和余指标的偶性及其拓扑后果
% * 无指标一和无余指标一情形的连通水平集论证
% 
%   主范围为紧连通 Hamilton 环面作用的矩映射分量；核验极值分量、连通性与偶指标条件，不把连通水平集结论宣称为任意 Morse–Bott 函数都成立。
% 
% 纲要内容结束

\section{矩映射的局部结构与 Morse–Bott 工具}
\label{g4:sec:1201}

矩映射的线性分量把环面作用转为实值函数。它的临界点正是相应圆方向停止运动的点；这些点往往组成子流形，普通 Morse 引理中的孤立临界点因而需要推广。本节先建立轨道附近的辛模型，再把法向二次型转成负盘丛的附着，最后证明本章所需的连通水平集结论。

\subsection{固定点附近的权模型与相容复法丛}
\label{g4:subsec:1201:01}

固定 \(T=\mathfrak t/(2\pi\Lambda)\)，其中 \(\Lambda\cong\mathbb Z^k\)。本章的 Hamilton 作用包含一个等变矩映射 \(\mu:M\to\mathfrak t^*\)；因为 \(T\) 交换，等变性就是 \(T\) 不变性。沿用正指数基本场及第十章的符号：
\[
	\xi_M(x)=\left.\frac{\mathrm d}{\mathrm dt}\right|_{t=0}
	\exp(t\xi)x,\qquad
	\mathrm d\mu^\xi=\iota_{\xi_M}\omega,
	\qquad \mu^\xi=\langle\mu,\xi\rangle.
\]
作用暂不要求有效。一个复直线上的权 \(\alpha\in\Lambda^*\) 按 \(\exp(\theta)z=e^{\mathrm i\langle\alpha,\theta\rangle}z\) 定义。若 \(z=x+\mathrm iy\)、\(\omega=\mathrm dx\wedge\mathrm dy\)，相应基本场为 \(\alpha(\xi)(-y\partial_x+x\partial_y)\)，所以其矩分量为 \(-\alpha(\xi)|z|^2/2\)。这个负号在全章保持不变。

\begin{proposition}\label{g4:prop:1201:invariant-complex-fixed-set}
	紧 Lie 群的辛作用具有不变的相容几乎复结构。若紧子群 \(K\) 作用于 \((M,\omega)\)，则固定点集 \(M^K\) 的每个连通分支是嵌入辛子流形，切空间为 \((\mathrm T_xM)^K\)。其辛法丛具有 \(K\) 不变的相容复结构。
\end{proposition}
\begin{proof}
	第一卷定理17.3.4用紧群平均给出不变正定度量 \(h\)。对 \((h,\omega)\) 应用\cref{g4:thm:0703:polar-construction}的逐点正规化，得到相容的 \(J\)；构造只用线性方程和正定平方根，与保持 \(h,\omega\) 的映射交换，故 \(J\) 不变。

	在固定点 \(x\) 附近，先取坐标 \(a\)，使 \(a(x)=0\)、\(\mathrm da_x=\operatorname{id}\)，并记 \(\rho(g)=\mathrm d(g\cdot)_x\)。紧性允许取一个共同邻域，使全部 \(a(gy)\) 都有定义。平均映射
	\[
		A(y)=\int_K\rho(g)^{-1}a(gy)\,\mathrm dg
	\]
	满足 \(A(gy)=\rho(g)A(y)\)、\(\mathrm dA_x=\operatorname{id}\)。逆函数定理给出等变局部坐标；把其像缩为一个不变小球，就得到 \(M^K\) 对应于线性固定子空间的坐标图。这证明嵌入性与切空间公式。固定子空间被 \(J\) 保持，而且 \(\omega(v,Jv)>0\) 对非零固定向量成立，故限制型非退化。其辛正交补也被 \(J\) 与 \(K\) 保持，给出所述复法丛。
\end{proof}

\begin{theorem}\label{g4:thm:1201:fixed-weight-model}
	设 \(T\) Hamilton 作用于实维数 \(2n\) 的无边界辛流形，\(x\in M^T\)。则存在以 \(x\) 为原点的等变辛坐标 \(z_j=x_j+\mathrm iy_j\) 和权 \(\alpha_1,\ldots,\alpha_n\in\Lambda^*\)，使
	\begin{equation}\label{g4:eq:1201:fixed-weight-model}
		\omega=\sum_{j=1}^n\mathrm dx_j\wedge\mathrm dy_j,
		\qquad
		\mu(z)=\mu(x)-\frac12\sum_{j=1}^n|z_j|^2\alpha_j.
	\end{equation}
	零权方向恰好组成固定子流形的切空间。
\end{theorem}
\begin{proof}
	取上一命题的不变相容复结构。切空间上的作用为交换的酉算子族。把一个非标量算子的特征子空间分开，其余算子保持这些子空间；对子空间维数归纳，便得到共同的酉特征直线。每条直线上的特征为连续字符 \(T\to S^1\)。沿 \(\mathfrak t\) 的一参数群，它写成 \(e^{\mathrm i\alpha(\theta)}\)；\(2\pi\Lambda\) 中的元素作用恒等，故 \(\alpha(\Lambda)\subset\mathbb Z\)。这给出整数权及标准辛线性坐标。

	上一命题的平均坐标将作用线性化，且拉回的辛形式 \(\omega_1\) 在原点等于常系数形式 \(\omega_0\)。用\cref{g4:thm:0803:relative-moser}对单点作相对变形。这里径向同伦与线性群作用交换，所以其原始一形式、收缩方程的解及流都不变或等变；在一个共同的不变小邻域内得到等变辛坐标。等变性与收缩方程说明，原矩映射和式\eqref{g4:eq:1201:fixed-weight-model}右边每个分量的微分相同。邻域连通，二者相差常数；在原点取值相同，便完全相等。零权结论来自线性固定子空间。\cite[第5.7节，定理5.22，第159页]{Cannas2006Symplectic}
\end{proof}

同一证明适用于任意紧交换群 \(H\) 的固定点，只须将整数权改为 \(H\) 的完整字符 \(\chi_j:H\to S^1\)。若 \(H\) 不连通，字符在有限分支上的取值不能由其微分 \(\alpha_j\in\mathfrak h^*\) 恢复。沿固定子流形移动时，各字符的子空间拼成复子丛：投影算子为 \(v\mapsto\int_H\overline{\chi(h)}\,h v\,\mathrm dh\)，它光滑且秩局部恒定。这也说明相容复法丛上的权分解具有丛意义。

\subsection{正则轨道与局部辛标准形}
\label{g4:subsec:1201:02}

在非固定轨道附近，还要保留轨道本身以及与轨道配对的横向变量。设 \(x\in M\)、\(H=T_x\)、\(\mathfrak h=\operatorname{Lie}(H)\)。因为矩映射沿轨道不变，\(E=\mathrm T_x(Tx)\) 各向同性。于是
\[
	V=E^\omega/E
\]
是带有 \(H\) 辛表示的向量空间，称为该轨道的\term[xinqiepian]{辛切片}[symplectic slice]。记 \(\mathfrak h^\circ\subset\mathfrak t^*\) 为消灭 \(\mathfrak h\) 的子空间，并选择分裂 \(\mathfrak t^*=\mathfrak h^\circ\oplus\mathfrak h^*\)。

\begin{theorem}\label{g4:thm:1201:orbit-normal-form}
	设 \(T\) Hamilton 作用于无边界辛流形 \((M^{2n},\omega)\)。对任意 \(x\in M\)，令 \(H=T_x\)，并取上述辛切片 \(V\)。轨道 \(Tx\) 有一个 \(T\) 等变辛邻域模型
	\begin{equation}\label{g4:eq:1201:orbit-model}
		Y=T\times_H(\mathfrak h^\circ\times V),
		\qquad
		(t,\eta,v)\sim(th^{-1},\eta,hv).
	\end{equation}
	零截面 \(T/H\) 对应于 \(Tx\)。给 \(V\) 取不变酉坐标，设完整字符为 \(\chi_1,\ldots,\chi_m\)，微分权为 \(\alpha_j\in\mathfrak h^*\)。模型的矩映射为
	\begin{equation}\label{g4:eq:1201:orbit-moment}
		\mu_Y([t,\eta,z])=\mu(x)+\eta+\nu_V(z),
		\qquad \nu_V(z)=-\frac12\sum_{j=1}^m|z_j|^2\alpha_j.
	\end{equation}
	此处保留整个稳定子 \(H\)，并有 \(\dim_{\mathbb R}V=2n-2\dim(T/H)\)。
\end{theorem}
\begin{proof}
	先构造模型的辛形式。在 \(\mathrm T^*T=T\times\mathfrak t^*\) 上取典范形式 \(\sum\mathrm d\theta_i\wedge\mathrm da_i\)。令 \(H\) 在 \(\mathrm T^*T\times V\) 上作用为
	\[
		h\cdot(t,a,v)=(th^{-1},a,hv).
	\]
	它的矩映射为 \(\nu_V(v)-a|_{\mathfrak h}\)：第一项来自切片，第二项来自负向平移。其微分在 \(a\) 方向满射，且整个 \(H\) 的作用自由，因为 \(th^{-1}=t\) 迫使 \(h=e\)。紧性给固有性，故\cref{g4:thm:1002:symplectic-reduction}适用，包括 \(H\) 不连通的情形。零水平满足 \(a=\eta+\nu_V(v)\)，所以约化商正是式\eqref{g4:eq:1201:orbit-model}。左 \(T\) 平移与约化相容，其矩映射为 \(a\)，再加常数 \(\mu(x)\) 得式\eqref{g4:eq:1201:orbit-moment}。在零截面处，模型切空间为
	\[
		(\mathfrak t/\mathfrak h)\oplus\mathfrak h^\circ\oplus V,
	\]
	辛形式在前两项之间为 \(\eta'(u)-\eta(u')\)，在第三项上为 \(\omega_V\)。

	接着把真实轨道的一阶数据与模型匹配。第一卷命题17.1.6给出 \(T/H\to M\) 的单射浸入；其源紧、目标 Hausdorff，故它还是到其像的同胚。因此轨道是紧嵌入子流形，可以使用管状邻域。取 \(H\) 不变相容复结构 \(J\)。\(E\oplus JE\) 为辛子空间，而 \((E\oplus JE)^\omega\) 规范地表示 \(E^\omega/E\)。\(H\) 在 \(E\) 上恒等，因为 \(T\) 交换；它也在 \(JE\) 上恒等。通过配对 \(w\mapsto(u\mapsto\omega(u,w))\)，将 \(JE\) 与 \(\mathfrak h^\circ\) 识别，便得到一个 \(H\) 等变辛线性同构，且在轨道切向为原有识别。用 \(T\) 平移将它延拓到整条轨道，得到两边环境切丛沿轨道的辛同构。

	第二卷定理15.3.5的管状邻域构造可在这里取为等变的。具体地，把构造中使用的度量和联络对 \(T\) 平均；联络的平均在联络这个仿射空间中进行，仍满足 Leibniz 规则。其局部指数由一个 \(T\) 不变微分方程给出，唯一性使指数等变。轨道紧，第二卷引理15.3.3的紧集缩小论证给统一的单射邻域。匹配上述法向映射后，得到固定轨道、且微分为指定辛同构的等变微分同胚 \(\Psi_0\)，从模型邻域映到真实邻域。

	因此 \(\Psi_0^*\omega\) 与模型形式沿轨道作为整个环境二张量相同。再次用\cref{g4:thm:0803:relative-moser}；若所取原始一形式还未不变，先对 \(T\) 平均，外微分及沿轨道的一阶消失条件都保持。收缩方程有唯一解，故变形流等变并逐点固定轨道。复合后得到所需等变辛映射。它保持矩映射的微分，而两矩映射在该轨道上取同一值，故在连通的小邻域上相等。
\end{proof}

当稳定子平凡时，\(\mathfrak h=0\)，模型成为 \(\mathrm T^*T\times V\)，矩映射只读出余切变量。若稳定子非平凡，横向的二次权项便产生矩像的边界或角。有限稳定子虽然不出现在 \(\mathfrak h\) 中，却仍作用于 \(V\) 并参与式\eqref{g4:eq:1201:orbit-model}的商；将 \(H\) 换成 \(H^0\) 一般会改变邻域。

\subsection{临界子流形、法向 Hessian 与 Morse–Bott 引理}
\label{g4:subsec:1201:03}

在临界点处，函数的 Hessian 与坐标无关，\footnote{Ludwig Otto Hesse（1811--1874），德国数学家，研究代数、分析与解析几何。}这是第二卷命题17.1.1的结论。若 \(C\) 是临界点组成的子流形，沿 \(C\) 微分恒为零，故再沿 \(C\) 求导得到 \(\mathrm T C\subset\ker\operatorname{Hess}f\)。因此 Hessian 自然下降到法丛 \(N_C=\mathrm TM|_C/\mathrm TC\)。

\begin{definition}\label{g4:def:1201:morse-bott}
	无边界光滑流形上的函数 \(f\) 称为\term[morsebotthanshu]{Morse–Bott 函数}[Morse--Bott function]，\footnote{Raoul Bott（1923--2005），匈牙利出生的美国数学家，研究拓扑学与微分几何。}若其临界点集是嵌入子流形的不交并，并且每个临界分支 \(C\) 都满足
	\[
		\ker\operatorname{Hess}_x f=\mathrm T_xC\qquad(x\in C).
	\]
	法向 Hessian 的负、正惯性指数分别称为该分支的指标和余指标，记作 \(\lambda_C,\nu_C\)。相应的负、正法丛记作 \(E_C^-,E_C^+\)。
\end{definition}

这些指标在连通分支上恒定，因为非退化对称型的惯性指数在小扰动下不变。选法丛度量，把 Hessian 看成可逆自伴算子 \(A\)；算子 \(A(A^2)^{-1/2}\) 的 \(\pm1\) 特征投影光滑，故确实给出两个光滑子丛。与第二卷普通 Morse 理论相比，新增的是临界点之间的切向方向及可能非平凡的 \(E_C^-\)，不能把所有负方向都当作同一个固定 Euclid 空间。

\begin{lemma}[Morse–Bott 引理]\label{g4:lem:1201:morse-bott}
	设 \(C\) 是 \(f\) 的一个紧连通临界子流形，\(f|_C=c\)，且法向 Hessian 非退化。则其一个邻域可与 \(E_C^-\oplus E_C^+\) 的零截面邻域微分同胚地识别，使
	\begin{equation}\label{g4:eq:1201:morse-bott-model}
		f(q,u,v)=c-\|u\|^2+\|v\|^2.
	\end{equation}
	范数来自两个丛上的光滑正定度量。
\end{lemma}
\begin{proof}
	用第二卷定理15.3.5取管状坐标 \((q,w)\)。沿每条法纤维作 Taylor 积分，得到
	\[
		f(q,w)-c=\langle A(q,w)w,w\rangle,
		\qquad
		A(q,w)=\int_0^1(1-t)\,\mathrm D_w^2f(q,tw)\,\mathrm dt.
	\]
	这里用一个固定丛度量识别对称型与算子，\(A_0(q)=A(q,0)\) 可逆。第二卷定理17.1.3逐变量配平方得到普通 Morse 坐标；为使坐标沿整个 \(C\) 相容，这里直接作一个丛算子的配平方。

	置 \(D(q,w)=A_0(q)^{-1}A(q,w)\)。紧性允许统一缩小，使 \(\|D-I\|<1\)。幂级数
	\[
		B=D^{1/2}=\sum_{m=0}^{\infty}\binom{1/2}{m}(D-I)^m
	\]
	在更小邻域上连同各阶导数一致收敛，定义光滑丛算子。它满足 \(B^2=D\)。由于 \(D^*A_0=A_0D\)，同一恒等式对其幂级数成立，故
	\[
		B^*A_0B=A_0B^2=A(q,w).
	\]
	所以变换 \((q,w)\mapsto(q,B(q,w)w)\) 把函数精确变成 \(c+\langle A_0(q)w,w\rangle\)。该变换沿零截面的微分为恒等，紧零截面上的局部逆与单射缩小论证使它在一个共同邻域内为微分同胚。最后按 \(A_0\) 的负、正子丛分解，并分别以 \(-\langle A_0u,u\rangle\)、\(\langle A_0v,v\rangle\) 为范数平方，即得式\eqref{g4:eq:1201:morse-bott-model}。
\end{proof}

\subsection{负法丛盘丛与跨临界值的次水平集变化}
\label{g4:subsec:1201:04}

记 \(M^a=\{f\leq a\}\)。第二卷命题17.2.1证明，在没有临界点的紧高度带上，向量场 \(\operatorname{grad}f/\|\operatorname{grad}f\|^2\) 的流使高度以速度一变化，因而给出乘积领带。跨过临界值时，需要把上一引理的二次模型接入这条流。

\begin{proposition}\label{g4:prop:1201:critical-crossing}
	设 \(f\) 为紧无边界流形上的 Morse–Bott 函数。若闭带 \(f^{-1}([c-\varepsilon,c+\varepsilon])\) 内的临界点恰为高度 \(c\) 上的连通分支 \(C\)，则充分缩小正则带宽后，\(M^{c+\varepsilon}\) 由 \(M^{c-\varepsilon}\) 附着
	\begin{equation}\label{g4:eq:1201:bundle-handle}
		D(E_C^-)\mathbin{\times_C}D(E_C^+)
		\quad\text{沿}\quad
		S(E_C^-)\mathbin{\times_C}D(E_C^+)
	\end{equation}
	并磨圆接缝得到。它相对旧区域同伦等价于附着负法丛盘丛 \(D(E_C^-)\)，附着边界为 \(S(E_C^-)\)。在正则水平集上，相应变化是把 \(S(E_C^-)\times_C D(E_C^+)\) 的内部换成 \(D(E_C^-)\times_C S(E_C^+)\)。同一临界值上有多个分支时，对互不相交的邻域同时作这些附着。
\end{proposition}
\begin{proof}
	在式\eqref{g4:eq:1201:morse-bott-model}的丛邻域中取 \(0<r^2<\varepsilon\)，使下列闭块仍位于该邻域：
	\[
		\mathcal H=\{(q,u,v):\|v\|\leq r,
		\ \|u\|^2\leq\varepsilon+\|v\|^2\}.
	\]
	映射
	\[
		(q,a,b)\longmapsto
		\bigl(q,\sqrt{\varepsilon+r^2\|b\|^2}\,a,rb\bigr)
	\]
	将单位盘丛乘积微分同胚地送到 \(\mathcal H\)。其与旧次水平集的交恰为 \(\|a\|=1\)，因为该处 \(f=c-\varepsilon\)；其余点的高度较大。因而已得到式\eqref{g4:eq:1201:bundle-handle}的附着面。

	丛上的竖直向量场 \(X=(2u,-2v)\) 满足
	\[
		Xf=-4(\|u\|^2+\|v\|^2),\qquad
		X(\|v\|^2-r^2)=-4r^2\quad(\|v\|=r).
	\]
	这些公式在正交丛过渡下不变，故是沿整个 \(C\) 的构造。把 \(X\) 与外部的负梯度场用单位分解拼接，保持它在 \(\mathcal H\) 附近的表达；拼接后在临界集外仍有 \(Xf<0\)。令 \(A=M^{c-\varepsilon}\)、\(K=A\cup\mathcal H\)。流在旧水平及新块自由面上都严格朝向 \(K\)。

	这里第二卷定理17.2.4的磨圆和流搬运可逐纤维进行，其具体横向函数仍为 \(s=f-c+\varepsilon\)、\(t=\|v\|^2-r^2\)。在接缝处将 \(\min(s,t)\) 替换为
	\[
		m_\delta(s,t)=\tfrac12\bigl(s+t-\eta_\delta(s-t)\bigr),
	\]
	其中 \(\eta_\delta\geq|\cdot|\)、\(|\eta_\delta'|\leq1\)，且离零足够远时等于绝对值。两偏导非负且和为一，故 \(Xm_\delta<0\)。磨圆后的 \(K_\delta\) 因而有被 \(X\) 严格穿入的光滑边界。剩下的紧带不含临界点，\(-Xf\) 有正下界；每条下降轨道在有限时间进入 \(K_\delta\)，反向则到达上水平。两端横截性使首次到达时间光滑，所以剩余部分是边界的一条领带。把该领带沿流线压入内部领圈，得到 \(M^{c+\varepsilon}\) 与磨圆附着空间的微分同胚。指标或余指标为零时，省去相应空的附着面或自由面即可。

	同伦结论也保留丛的扭转。第二卷引理17.2.5的核心收缩只依赖纤维范数：在单位盘坐标 \((a,b)\) 中置 \(T(a,b)=\min\{\|b\|,-\log\|a\|\}\)，约定 \(T(0,b)=\|b\|\)，并用
	\[
		(a,b)\longmapsto
		\left(e^{sT}a,\left(1-\frac{sT}{\|b\|}\right)b\right),
		\qquad0\leq s\leq1,
	\]
	其中 \(b=0\) 时取恒等。该公式连续，逐点固定 \(\|a\|=1\) 的附着面，最终落入附着面与 \(b=0\) 核心的并。它与正交过渡相容，故在整个盘丛上成立，并与旧区域的恒等映射粘合。最后，盘丛乘积的边界由式\eqref{g4:eq:1201:bundle-handle}的两个面组成；移去附着面，留下的正是 \(D(E_C^-)\times_C S(E_C^+)\)，得到所述水平集变化。多个临界分支的邻域可取互不相交，全部构造遂可同时进行。
\end{proof}

\subsection{一般圆方向矩映射分量的 Morse–Bott 性}
\label{g4:subsec:1201:05}

\begin{proposition}\label{g4:prop:1201:moment-morse-bott}
	设紧环面 \(T\) Hamilton 作用于紧无边界辛流形 \(M\)。对任意 \(\xi\in\mathfrak t\)，函数 \(f=\mu^\xi\) 是 Morse–Bott 函数。若 \(S=\overline{\exp(\mathbb R\xi)}\)，则
	\[
		\operatorname{Crit}(f)=M^S.
	\]
	沿其每个连通分支 \(C\)，负、正法丛都是复向量丛，因而指标与余指标都是偶数。
\end{proposition}
\begin{proof}
	非退化性给出 \(\mathrm df=0\) 当且仅当 \(\xi_M=0\)，又由流的唯一性，这等价于整个一参数群固定该点；连续性把固定性延伸到闭包 \(S\)。第一卷引理19.2.1及定理19.2.3说明，\(S\) 是紧连通交换 Lie 群，因而是子环面。\cref{g4:prop:1201:invariant-complex-fixed-set}说明其固定集是辛子流形。

	在 \(x\in C\) 处对 \(S\) 使用固定点权模型。零权是 \(\mathrm T_xC\)，其余权记为 \(\beta_j\)。若 \(\beta_j(\xi)=0\)，该字符在稠密的一参数子群上恒等，因而在 \(S\) 上恒等，与它是非零权矛盾。因此
	\[
		f=f(x)-\frac12\sum_j\beta_j(\xi)(x_j^2+y_j^2)
	\]
	在全部法向上具有非退化 Hessian，核恰为 \(\mathrm T_xC\)。\(\beta_j(\xi)>0\) 的复方向为负方向，\(\beta_j(\xi)<0\) 的复方向为正方向。权子丛的光滑拼接给出复负、正法丛，并有
	\begin{equation}\label{g4:eq:1201:even-indices}
		\lambda_C=2\sum_{\beta(\xi)>0}\operatorname{rank}_{\mathbb C}E_\beta,
		\qquad
		\nu_C=2\sum_{\beta(\xi)<0}\operatorname{rank}_{\mathbb C}E_\beta.
	\end{equation}
	固定集紧，每点有只包含一个局部固定片的坐标邻域；有限覆盖说明它只有有限个连通分支，故也只有有限个临界值。
\end{proof}

“圆方向”通常指 \(\xi\in\Lambda\)，此时 \(\exp(t\xi)\) 以 \(2\pi\) 为一个周期；上述结论也覆盖非周期的一参数方向。当 \(k\geq1\) 时，还可以选一个整数方向，使它的临界集恰为 \(M^T\)。事实上，轨道模型中出现的稳定子 Lie 代数是 \(\mathfrak h\) 内若干权核的交，故每张模型只有有限种；紧性给有限张模型覆盖，因此全局也只有有限种。选 \(\xi\in\Lambda\) 避开其中所有真子空间即可。这样的整数向量存在：有限个非零线性方程不能消灭全部 \((1,N,\ldots,N^{k-1})\)，每个方程只排除有限多个整数 \(N\)。

\begin{example}\label{g4:exm:1201:opposite-weights}
	在 \(\mathbb C^2\) 上让圆以权 \((1,-1)\) 作用。矩函数为 \(f=(|z_2|^2-|z_1|^2)/2\)。原点的指标、余指标都是 \(2\)，其零水平集是锥 \(\{|z_1|=|z_2|\}\)，并非光滑超曲面。偶指标控制的是跨临界高度的拓扑变化，并不使临界水平本身变成正则水平。
\end{example}

\subsection{指标和余指标的偶性及其拓扑后果}
\label{g4:subsec:1201:06}

对于指标 \(\lambda_C\geq2\) 的连通临界分支，附着面 \(S(E_C^-)\times_C D(E_C^+)\) 连通，因为底空间连通且球面纤维 \(S^{\lambda_C-1}\) 连通。因此这个附着面只能落在旧次水平集的一个连通分支中；附着不能把两个旧分支合并。指标为零的分支则没有附着面，会新生一个连通分支。指标一恰是可能合并分支的情形。

\begin{proposition}\label{g4:prop:1201:unique-extreme-components}
	设 \(f\) 是紧连通无边界流形上的 Morse–Bott 函数，临界分支有限。若没有指标一的临界分支，则它只有一个指标零的临界分支，该分支正是全局最小值集。若同时没有余指标一的临界分支，则最大值集也只有一个连通分支，且每个非极值临界分支的指标、余指标都至少为二。
\end{proposition}
\begin{proof}
	从低于最小值的空次水平集开始，依次越过有限个临界值。无临界带保持拓扑；\cref{g4:prop:1201:critical-crossing}及前面的连通性说明，指标零分支各新生一个组件，而全部正指标分支都不能合并旧组件。最高次水平集是连通的整个流形，因此总共只能新生一个组件。全局最小值必在指标零的临界分支上，故它就是唯一的这一分支；任何局部极小临界分支也必须等于它。对 \(-f\) 重复论证，得到最大值结论。其余分支既不是指标零也不是余指标零，再排除一，即得最后的断言。若 \(f\) 常值，唯一分支就是整个流形，结论中的非极值部分为空。
\end{proof}

矩映射分量满足该命题的两组假设。特别地，它不会在一个中间高度新生孤立的局部极小组件；下降和上升两个方向的偶性都在这里发挥作用。

\subsection{无指标一和无余指标一情形的连通水平集论证}
\label{g4:subsec:1201:07}

\begin{theorem}\label{g4:thm:1201:component-connected-levels}
	紧连通无边界辛流形上，Hamilton 环面作用的每个矩映射分量 \(\mu^\xi\) 的非空水平集都连通，包括临界水平。更一般地，同样结论适用于紧连通无边界流形上临界分支有限、且既无指标一也无余指标一分支的 Morse–Bott 函数。
\end{theorem}
\begin{proof}
	先证更一般的陈述。常值情形直接成立。否则最小值集是上一命题给出的唯一连通临界分支 \(C_{\min}\)，其余指标 \(\nu\) 至少为二。在最小值上方足够近的正则水平，Morse–Bott 模型给出 \(S(E_{C_{\min}}^+)\)，它连通。

	设已知某个中间临界值下方的正则水平连通。此高度上的每个临界分支都满足 \(\lambda,\nu\geq2\)。由\cref{g4:prop:1201:critical-crossing}，向上穿越时从旧水平移去 \(S(E^-)\times_C D(E^+)\) 的内部，再贴上 \(D(E^-)\times_C S(E^+)\)。旧水平中的核心 \(S(E^-)\) 余维为 \(\nu\geq2\)，移去其盘邻域的内部不会破坏连通性。为说明这一点，在适配坐标中，余维至少二的线性子空间的补集道路连通；因此子流形各点都有补集连通的小邻域。若整个补集分为不同组件，这些邻域会把每个被删点唯一分配到其中一个组件，使原流形也分为不交开集，矛盾。再沿法纤维径向推出盘邻域，就得到所需结论。新贴入的 \(D(E^-)\times_C S(E^+)\) 也连通，且通过非空的 \(S(E^-)\times_C S(E^+)\) 接到剩余部分。因此新正则水平仍连通。对互不相交的全部临界分支同时进行，并在无临界带中搬运，归纳得到所有正则水平连通。

	最后处理一个中间临界值 \(c\)。每个临界分支都有负方向，所以式\eqref{g4:eq:1201:morse-bott-model}表明其每一点都可由高度 \(c-\delta\) 的点逼近，\(\delta\downarrow0\)。水平 \(f=c\) 上的正则点也可用局部高度坐标这样逼近。另一方面，紧性保证高度趋于 \(c\) 的任意点列的聚点都位于 \(f^{-1}(c)\)。若该紧水平集分成两个非空不交闭集，可取它们的不交开邻域；充分接近 \(c\) 的下方正则水平必包含在两邻域的并中，并且分别与二者相交。这与这些正则水平连通矛盾。最小、最大水平已经由上一命题处理。

	对 \(f=\mu^\xi\)，\cref{g4:prop:1201:moment-morse-bott}给出所有假设，特别是正负法丛的偶实秩。因此得到矩映射分量的结论。
\end{proof}

上述证明同时使用了紧性、连通性和两侧指标限制。普通 Morse–Bott 函数若有指标一或余指标一的临界分支，水平集可以断开，例如标准二维环面的高度函数。下一节将进一步处理若干分量同时固定的纤维；单独知道每个实值分量的水平集连通，还需要额外的局部到整体论证。
% 纲要条目排为编号小节标题，并进入目录与书签；说明段仍保留于纲要注释中。

% !TeX root = ../../Book4.tex
% 纲要标题：### 12.2 凸性定理
% 纲要位置：几何学计划纲要.md，第 3933 行。
% 正文状态：初稿；尚未正式审核。

% 纲要内容开始
% 
% * 圆方向的极值和连通水平集
% * 环面秩的归纳
% * 矩映射像的凸性
% * 固定点像及凸包描述
% 
% 纲要内容结束

\section{凸性定理}
\label{g4:sec:1202}

本节证明紧连通 Hamilton 环面流形的全部矩纤维连通，并证明矩像是固定点像的凸包。两个步骤各有职责：环面秩的归纳在极值分支中找到共同固定点；轨道局部模型则把纤维连通分支组成的商变成局部凸空间。这个紧商上的最短路将同时给出纤维连通与像的凸性。

\subsection{圆方向的极值和连通水平集}
\label{g4:subsec:1202:01}

以下 \(M\) 为紧连通无边界辛流形，\(T\) 为紧环面，\(\mu:M\to\mathfrak t^*\) 为等变矩映射。对任意 \(\xi\in\mathfrak t\)，\cref{g4:thm:1201:component-connected-levels}说明
\[
	\{x\in M:\langle\mu(x),\xi\rangle=c\}
\]
非空时连通。特别地，\(\mu^\xi\) 的最小值集与最大值集各为一个连通辛子流形。若 \(T=S^1\)，矩映射只有一个分量，这已经证明了所有矩纤维连通；其像由紧性与连通性成为闭区间。

\ProseParagraphBreak

当 \(\dim T>1\) 时，固定一个线性分量得到的是矩像中一个超平面的原像，而完整矩纤维要求所有分量同时固定。后面的论证将保留这些交集的信息。仅仅知道所有线性投影的像都是区间，甚至不能排除一个平面圆周作为候选像；因此投影区间性质本身不能证明高维凸性。

\subsection{环面秩的归纳}
\label{g4:subsec:1202:02}

\begin{proposition}\label{g4:prop:1202:fixed-point-induction}
	任意紧连通无边界 Hamilton 环面流形都有共同固定点。更具体地，若 \(S\subset T\) 是正维子环面，则 \(M^S\) 的每个非空连通分支都包含 \(T\) 固定点。因此任意矩映射分量的每个极值分支都包含 \(T\) 固定点。
\end{proposition}
\begin{proof}
	先按 \(k=\dim T\) 归纳证明共同固定点存在。\(k=0\) 时每一点都固定。设结论已对更低维环面成立，在 \(T\) 中取一个圆子群 \(S\)，令 \(\xi\) 为它的非零生成元。函数 \(\mu^\xi\) 在紧空间上达到最大值，取最大值集中的一个连通分支 \(C\)。上一节说明 \(C\) 是紧连通辛子流形，且 \(S\) 逐点固定它。\(T\) 保持 \(C\)：作用保持高度，又因 \(T\) 连通，不能把一个点连续移到固定集的另一个连通分支。

	\(T/S\) 在 \(C\) 上的作用仍是 Hamilton 的。对每个 \(\zeta\in\mathfrak s\)，有 \(\zeta_M|_C=0\)，所以 \(\mu^\zeta|_C\) 为常数。选 \(a\in\mathfrak t^*\) 与 \(\mu|_C\) 在 \(\mathfrak s\) 上相同，则
	\[
		\nu=\mu|_C-a:C\longrightarrow
		\mathfrak s^\circ\cong(\mathfrak t/\mathfrak s)^*
	\]
	满足商环面所需的收缩方程及不变性。\(\dim(T/S)=k-1\)，归纳假设给出它在 \(C\) 上的固定点，而该点也被 \(S\) 固定，故被整个 \(T\) 固定。

	现在令 \(S\) 为任意正维子环面，\(C\) 为 \(M^S\) 的一个连通分支。固定集的辛性、紧性和 \(T\) 不变性仍成立，同样构造商环面的矩映射。刚刚完成的归纳结果用于 \(T/S\) 作用于 \(C\)，便得到 \(T\) 固定点。对任意非零 \(\xi\)，\(\mu^\xi\) 的极值分支是 \(M^{\overline{\exp(\mathbb R\xi)}}\) 的分支，由此得到最后结论；\(\xi=0\) 时极值集为整个 \(M\)，使用第一部分即可。
\end{proof}

这一归纳没有要求作用有效。即使某个圆子群在整个流形上作用恒等，商环面的维数仍降低，归纳仍然有效。对任意 \(\xi\)，我们已经得到
\begin{equation}\label{g4:eq:1202:fixed-support-extreme}
	\max_{x\in M}\langle\mu(x),\xi\rangle
	=\max_{x\in M^T}\langle\mu(x),\xi\rangle.
\end{equation}
它将在凸性成立后确定整个多面体。

\subsection{矩映射像的凸性}
\label{g4:subsec:1202:03}

先从轨道模型提取精确的局部信息。一个由有限向量 \(v_1,\ldots,v_N\) 的非负线性组合组成的集合称为有限生成凸锥。这里允许锥含直线，也允许其维数小于环境空间维数；所有“相对开”都指在该像自身的子空间拓扑中开。

\begin{lemma}\label{g4:lem:1202:linear-cone-openness}
	设 \(C=\mathbb R^d\times[0,\infty)^m\)，\(A:\mathbb R^{d+m}\to E\) 为线性映射。则 \(A(C)\) 是闭的有限生成凸锥，且 \(A:C\to A(C)\) 是开映射。
\end{lemma}
\begin{proof}
	\(C\) 由有限个向量生成：自由方向取正、负标准基，非负方向只取正标准基。其像因此有限生成。对任何非负表示 \(y=\sum c_jv_j\)，若正系数所对应的 \(v_j\) 线性相关，取一个非零关系 \(\sum a_jv_j=0\)，改变符号使某个 \(a_j>0\)，再减去 \((\min_{a_j>0}c_j/a_j)(a_j)\)。所有新系数仍非负，至少一个正系数变为零。有限步后得到线性无关的生成子集。每个这种子集上的逆线性映射有有界算子范数；子集只有有限种，所以可统一控制系数范数为 \(K\|y\|\)。这说明有界像点列可选有界系数表示，取子列即可证明锥闭。

	固定 \(x\in C\)。令 \(D_x\) 为如下方向锥：在 \(x_j=0\) 的非负坐标上只许非负增量，在 \(x_j>0\) 的坐标及自由坐标上允许正负增量。它也有限生成。对 \(y\in A(C)\)，任取 \(z\in C\) 使 \(Az=y\)，则 \(z-x\in D_x\)，故 \(y-Ax\in A(D_x)\)。把第一段的有界系数论证用于 \(A(D_x)\)，可选 \(u\in D_x\)，满足
	\[
		Au=y-Ax,\qquad \|u\|\leq K_x\|y-Ax\|.
	\]
	当右边足够小时，\(x\) 的正坐标不会因加上 \(u\) 而变成负数；原来的零坐标也保持非负。因此 \(x+u\in C\)，且可落在 \(x\) 的任意给定相对邻域内。这正说明该邻域的像包含 \(Ax\) 在 \(A(C)\) 中的一个邻域，故映射开。
\end{proof}

\begin{lemma}\label{g4:lem:1202:local-convex-fibres}
	每条 \(T\) 轨道 \(Tx\) 都有任意小的 \(T\) 不变邻域 \(U\)，使 \(\mu(U)\) 凸，\(\mu|_U\) 对其像为开映射，且 \(U\cap\mu^{-1}(b)\) 非空时道路连通。还可使
	\begin{equation}\label{g4:eq:1202:local-moment-cone}
		\mu(U)=\bigl(\mu(x)+C_x\bigr)\cap B_\varepsilon(\mu(x)),
		\qquad
		C_x=\mathfrak h^\circ+
		\operatorname{cone}(-\alpha_1,\ldots,-\alpha_m),
	\end{equation}
	其中 \(H=T_x\)，\(\alpha_j\) 是轨道模型中的切片微分权。
\end{lemma}
\begin{proof}
	在\cref{g4:thm:1201:orbit-normal-form}中先取 \(\eta\) 的小凸球和 \(V\) 的小球，得到 \(U=T\times_H(B\times B_V)\)。令 \(r_j=|z_j|^2/2\)，其像写成
	\[
		\mu(x)+B-\left\{\sum_jr_j\alpha_j:
		r_j\geq0,\ \sum_jr_j<\delta\right\},
	\]
	是凸集。映射 \(z\mapsto(r_j)\) 对非负正交象限开：非零复坐标的半径可以在开区间内变化，零坐标的小圆盘则给出相对区间 \([0,\epsilon)\)。上一引理再说明 \((\eta,r)\mapsto\eta-\sum r_j\alpha_j\) 对其锥像开。结合商映射的开性，得到 \(\mu|_U\) 的相对开性。

	取同一纤维中的两点，选代表 \((t_0,\eta_0,z^0)\)、\((t_1,\eta_1,z^1)\)。在 \(T\) 中连接 \(t_0,t_1\)；对 \(\eta,r\) 作线性插值。约束 \(\eta-\sum r_j\alpha_j=b-\mu(x)\) 及所选球、单纯形限制都保持。再写 \(z_j^a=\sqrt{2r_j^a}e^{\mathrm i\phi_j^a}\)，零坐标的相位任取，令
	\[
		z_j(s)=\sqrt{2((1-s)r_j^0+sr_j^1)}
		e^{\mathrm i((1-s)\phi_j^0+s\phi_j^1)}.
	\]
	这给出一条连续纤维内道路，下降到 \(H\) 商后仍连接原两点。

	上一引理的有界表示还说明，\(C_x\) 中充分小的向量都可用充分小的 \(\eta,r\) 表示。因此该凸像在 \(\mu(x)\) 附近与 \(\mu(x)+C_x\) 相同。进一步把 \(U\) 与 \(\mu^{-1}(B_\varepsilon(\mu(x)))\) 相交，就得到式\eqref{g4:eq:1202:local-moment-cone}；每条纤维的原片没有变化，相对开性也保持。
\end{proof}

把同一矩纤维的每个连通分支压成一点，记所得商空间为 \(Q\)，商映射为 \(q:M\to Q\)。矩映射分解为 \(\mu=\overline\mu\circ q\)。这时尚未知道一条矩纤维有几个分支，所以 \(\overline\mu\) 暂时可能多对一。下面先在商上构造局部图。

\begin{lemma}\label{g4:lem:1202:fibre-component-quotient}
	\(Q\) 是紧连通 Hausdorff 空间。每个 \(z\in Q\) 都有开邻域 \(W\)，使 \(\overline\mu|_W\) 为到一个闭有限生成凸锥与开球之交的同胚；锥的顶点为 \(\overline\mu(z)\)。
\end{lemma}
\begin{proof}
	固定 \(a\in\mu(M)\)。上一引理给出 \(\mu^{-1}(a)\) 的一个有限开覆盖，每个局部纤维片连通，所以该纤维只有有限个连通分支 \(C_1,\ldots,C_s\)。它们紧且彼此不交，可放入彼此不交的开邻域中。对每个 \(C_j\)，在该开邻域内选有限个轨道模型 \(U_{ji}\) 覆盖它，且每个模型的中心在 \(C_j\)。

	对同一个 \(j\)，以 \(U_{ji}\cap C_j\) 为顶点、以在 \(C_j\) 中相交为边，所得有限交图连通，否则这些开片会把 \(C_j\) 分开。若两片在纤维上相交，\(\mu(U_{ji}\cap U_{j\ell})\) 包含 \(a\)，并且在两片的矩像中都相对开。因此两片的像在 \(a\) 附近具有相同的锥胚：在足够小的球内，每个像都包含于这个共同开像，故两者相等。沿交图传播，全部 \(U_{ji}\) 具有同一个锥胚 \(a+C_j'\)。有限性允许统一缩小球半径，使每片的像及交图每条边的交像都包含 \((a+C_j')\cap B_\varepsilon(a)\)，并使每片的像在此球内恰为这个集合。

	还可缩小 \(\varepsilon\)，使
	\[
		\mu^{-1}(B_\varepsilon(a))\subset\bigcup_{j,i}U_{ji}.
	\]
	否则可在右边补集内取矩值趋于 \(a\) 的点列，紧性给出落在 \(\mu^{-1}(a)\) 中的聚点，与覆盖矛盾。令 \(V_j=(\bigcup_iU_{ji})\cap\mu^{-1}(B_\varepsilon(a))\)。各 \(V_j\) 彼此不交，并在这个整原像中同时开闭。对每个 \(b\in(a+C_j')\cap B_\varepsilon(a)\)，各片的 \(b\) 纤维连通，且交图的每条边仍有非空纤维交。因此 \(V_j\cap\mu^{-1}(b)\) 连通，恰为整条矩纤维的一个连通分支。

	所以 \(V_j\) 是由完整等价类组成的开集，\(q(V_j)\) 在 \(Q\) 中开。局部相对开性还说明 \(\overline\mu:q(V_j)\to(a+C_j')\cap B_\varepsilon(a)\) 是开连续双射，故是同胚。由此得到所需图。若 \(Q\) 中两点矩值不同，用不交的矩值球分开；若矩值相同但对应分支不同，用上述不交的 \(q(V_j)\) 分开。因此 \(Q\) Hausdorff。紧性与连通性则来自它是紧连通空间 \(M\) 的连续商。
\end{proof}

\begin{theorem}\label{g4:thm:1202:connected-fibres}
	紧连通无边界辛流形上 Hamilton 环面作用的每条非空矩纤维 \(\mu^{-1}(a)\) 都连通，而且 \(\mu(M)\) 是凸集。
\end{theorem}
\begin{proof}
	利用上一引理的凸图，在 \(Q\) 上定义路径的长度为其 \(\overline\mu\) 像的 Euclid 长度，并令
	\[
		d(z,w)=\inf_\gamma\operatorname{length}(\overline\mu\circ\gamma),
	\]
	其中 \(\gamma\) 连接 \(z,w\)。可只取有限条凸图内直线段的拼接。这样的路径连接全部点：从任一点可达的集合开，其补集也开，\(Q\) 连通使可达集为整个 \(Q\)。所以 \(d\) 有限，且满足对称性及三角不等式。

	需要证明 \(d\) 真正分开点。取以 \(z\) 为中心的一个锥图 \(W_R\)，像为 \((a+C)\cap B_R(a)\)，其中 \(a=\overline\mu(z)\)。若 \(0<r<R\)，较小闭锥球的逆像紧，因 \(Q\) Hausdorff，它在 \(Q\) 中闭且包含于 \(W_R\)。因此一条从 \(z\) 出发、离开较小开锥球 \(W_r\) 的道路必须先达到 \(\|\overline\mu-a\|=r\)，长度至少为 \(r\)。对 \(w\in W_r\)，凸图中的直线段给出
	\[
		d(z,w)=\|\overline\mu(z)-\overline\mu(w)\|,
	\]
	因为任意像道路的长度都不小于端点距离。因此图内不同点被分开，图外点也与 \(z\) 至少相距 \(r\)。这证明 \(d\) 为度量，并且小度量球恰是小锥球的逆像，故度量拓扑与原拓扑一致。\(Q\) 因而是紧度量空间。在任一凸图内，两点间的 \(d\) 距离都等于其像点的 Euclid 距离。

	再证最短路存在。取连接 \(z,w\) 的折线路径 \(\gamma_j\)，长度 \(L_j\downarrow d(z,w)\)，按弧长在 \([0,1]\) 上重新参数化。它们关于 \(d\) 的 Lipschitz 常数可统一控制。紧度量空间中，先在可数稠密时间集上逐次取收敛子列，再用统一 Lipschitz 界，就得到一个在整个区间上一致收敛的子列及极限道路 \(\gamma\)。对每个有限时间分割，距离和连续地取极限，因此极限长度不大于 \(\lim L_j\)。局部图是 Euclid 等距图，故用 \(d\) 计算的曲线长度与像曲线长度相同。于是 \(\gamma\) 的长度恰为 \(d(z,w)\)，是最短路。

	记 \(D=d(z,w)\)。极限道路为 \(D\)-Lipschitz；由端点距离和三角不等式，对任意 \(s<t\) 还有
	\[
		D\leq Ds+d(\gamma(s),\gamma(t))+D(1-t),
	\]
	所以 \(d(\gamma(s),\gamma(t))=D(t-s)\)。在足够小的时间段内，道路落在一个凸图中；Euclid 三角不等式取等说明其像是一条匀速直线段。相邻时间段在开区间上重叠，速度向量因而相同。故整个像是直线段
	\[
		\overline\mu(\gamma(t))
		=(1-t)\overline\mu(z)+t\overline\mu(w),
		\qquad
		D=\|\overline\mu(z)-\overline\mu(w)\|.
	\]
	若两端矩值相同，就有 \(D=0\)，从而 \(z=w\)。因此 \(\overline\mu\) 单射，每条矩纤维只有一个连通分支，证明纤维连通。对任意两个矩值，上述最短路的像包含它们之间的整条线段，证明凸性。
\end{proof}

这个证明也解释了紧性与局部模型各自的作用。紧性使纤维只有有限个局部分支，并保证长度最小值可达到；局部凸锥及局部连通纤维使最短路的像不能在一个角处折返。最后得到的是完整矩纤维的连通性，包含秩下降的点。

\subsection{固定点像及凸包描述}
\label{g4:subsec:1202:04}

固定点集 \(M^T\) 是紧子流形并具有有限个连通分支。在每个分支上，全部 \(\xi_M\) 为零，故全部 \(\mathrm d\mu^\xi\) 为零；因此矩映射在该分支上常值。记这些有限个像为 \(a_1,\ldots,a_N\)。

下面的结论称为 Atiyah–Guillemin–Sternberg 凸性定理。\footnote{Michael Francis Atiyah（1929--2019），英国数学家，研究拓扑 K 理论与椭圆算子指标。Victor William Guillemin（1937年生），美国数学家，研究微局部分析与辛几何。Shlomo Zvi Sternberg（1936--2024），美国数学家，研究微分几何、辛几何与 Lie 理论。}

\begin{theorem}[Atiyah–Guillemin–Sternberg 凸性定理]\label{g4:thm:1202:convexity}
	设紧环面 \(T\) Hamilton 作用于紧连通无边界辛流形 \((M,\omega)\)，等变矩映射为 \(\mu\)。则所有非空矩纤维连通，且
	\begin{equation}\label{g4:eq:1202:fixed-convex-hull}
		\mu(M)=\operatorname{conv}\mu(M^T)
		=\operatorname{conv}\{a_1,\ldots,a_N\}.
	\end{equation}
	特别地，矩像是紧凸多面体，每个顶点都是固定点的矩像；作用不必有效，多面体也可以处于 \(\mathfrak t^*\) 的真仿射子空间中。
\end{theorem}
\begin{proof}
	纤维连通与凸性已经由\cref{g4:thm:1202:connected-fibres}证明。令 \(P=\operatorname{conv}\{a_1,\ldots,a_N\}\)。它是紧单纯形的连续仿射像，故紧；\(a_i\in\mu(M)\) 与矩像凸性给出 \(P\subset\mu(M)\)。

	若有 \(b\in\mu(M)\setminus P\)，在 \(P\) 中选一个距 \(b\) 最近的点 \(p\)，令 \(\xi=b-p\)，用 Euclid 内积将它视为 \(\mathfrak t\) 的元素。对任意 \(y\in P\)，线段 \(p+t(y-p)\) 位于 \(P\) 中。最小距离在 \(t=0\) 的单侧导数非负，故
	\[
		\langle\xi,y-p\rangle\leq0,
		\qquad
		\langle\xi,b\rangle
		=\langle\xi,p\rangle+\|b-p\|^2
		>\max_{y\in P}\langle\xi,y\rangle.
	\]
	这与式\eqref{g4:eq:1202:fixed-support-extreme}矛盾。因此 \(\mu(M)=P\)。有限集合的凸包中，一个极端点只能是生成点之一，否则它的凸组合表达会把它写成不同两点的真凸组合。所以每个顶点来自固定点。\cite[第5.7节，定理5.21，第158页]{Cannas2006Symplectic}
\end{proof}

式\eqref{g4:eq:1202:fixed-convex-hull}中的集合称为\term[juduomianti]{矩多面体}[moment polytope]。一个固定分支的像可以位于多面体内部，因而不能把所有固定点都与顶点一一对应。

\begin{example}\label{g4:exm:1202:diagonal-spheres}
	在两个单位球面上分别取面积形式 \(\omega_j=\mathrm d\theta_j\wedge\mathrm dz_j\)，其中 \(z_j\) 为高度；该表达在极点外成立，并光滑延拓为球面面积形式。对绕竖轴的正旋转，\(\iota_{\partial_{\theta_j}}\omega_j=\mathrm dz_j\)，所以高度就是矩函数。乘积环面作用于 \(S^2\times S^2\) 的矩像为正方形 \([-1,1]^2\)。限制到对角圆后，矩函数变为 \(z_1+z_2\)，像为 \([-2,2]\)。四个固定点的矩值依次为 \(-2,0,0,2\)，中间两个固定点都映到区间内部。凸包公式同时正确描述这两种不同秩的作用。
\end{example}

从局部权模型到有限凸包，我们已经得到一般紧 Hamilton 环面作用的整体约束。后面再加上有效性及半维条件，切片中的权数与稳定子维数相等，局部锥才会进一步满足 Delzant 的整数格条件。
% 纲要条目排为编号小节标题，并进入目录与书签；说明段仍保留于纲要注释中。

% !TeX root = ../../Book4.tex
% 纲要标题：### 12.3 Delzant 多面体与构造
% 纲要位置：几何学计划纲要.md，第 3940 行。
% 正文状态：初稿；尚未正式审核。

% 纲要内容开始
% 
% * 简单、有理和光滑格条件
% * 从多面体构造辛约化
% * 自由性和紧性的核验
% * 矩映射像的恢复
% 
% 纲要内容结束

\section{Delzant 多面体与构造}
\label{g4:sec:1203}

凸性定理从一个 Hamilton 环面作用得到多面体。本节反过来从多面体构造辛流形：给每个余维一面引入一个复坐标，把距离该面的线性函数变成这个坐标模长的平方，再用辛约化消去多余相位。哪些相位可以自由消去，正由顶点处的整数格条件决定。

\subsection{简单、有理和光滑格条件}
\label{g4:subsec:1203:01}

固定
\[
	N=\mathbb Z^n,\qquad
	\mathfrak t=N\otimes_{\mathbb Z}\mathbb R,\qquad
	T=\mathfrak t/(2\pi N),\qquad
	M_{\mathbb Z}=\operatorname{Hom}(N,\mathbb Z).
\]
多面体位于对偶空间 \(\mathfrak t^*=M_{\mathbb Z}\otimes\mathbb R\) 中。以下多面体均紧、凸、具有非空内部，且每个半空间不等式都对应一个实际的余维一面，不保留冗余不等式。

\begin{definition}\label{g4:def:1203:delzant-polytope}
	对一个 \(n\) 维凸多面体 \(P\subset\mathfrak t^*\)，若每个顶点恰落在 \(n\) 个余维一面上，则称它为\term[jiandanduomianti]{简单多面体}[simple polytope]。若各余维一面有整数法向量，则称它为关于 \(N\) 的\term[youliduomianti]{有理多面体}[rational polytope]；将法向量取为朝外的原始格向量后，可唯一写成
	\begin{equation}\label{g4:eq:1203:facet-presentation}
		P=\{x\in\mathfrak t^*:\langle x,\nu_j\rangle\le\lambda_j,\
		j=1,\ldots,d\},\qquad \nu_j\in N.
	\end{equation}
	这里“原始”表示 \(\nu_j\) 不是另一格向量的整数倍 \(k\nu\)、\(|k|>1\)。

	若 \(P\) 简单、有理，且每个顶点 \(v\) 所属的 \(n\) 个面的原始外法向量构成 \(N\) 的整数基底，则称 \(P\) 为\term[Delzantduomianti]{Delzant 多面体}[Delzant polytope]，这一基底要求称为光滑格条件。
\end{definition}

“有理”约束的是面的方向，并未要求顶点坐标为有理数或整数；支撑常数 \(\lambda_j\) 可以是任意实数。固定一个顶点，记相交面的指标集为 \(I_v=\{j_1,\ldots,j_n\}\)。若 \(m_1,\ldots,m_n\) 是 \(\nu_{j_1},\ldots,\nu_{j_n}\) 的对偶基，则从 \(v\) 出发的边沿 \(-m_i\) 方向。因此光滑格条件也等价于：每个顶点的原始内向边向量构成对偶格 \(M_{\mathbb Z}\) 的基底。负号来自我们选择了外法向量。

\ProseParagraphBreak

在二维，简单性对非退化凸多边形自动成立，光滑格条件则要求相邻两个原始外法向量组成的矩阵行列式为 \(\pm1\)。例如
\[
	\{x_1,x_2\ge0,\ x_1+x_2\le s\},\qquad
	[0,a]\times[0,b],\qquad s,a,b>0
\]
都满足条件。在第一个三角形中再加入 \(x_1+x_2\ge\varepsilon\)、\(0<\varepsilon<s\)，所得截角四边形也满足条件；新的原始外法向量是 \((-1,-1)\)。这些多面体分别预示射影空间、乘积和吹起的构造。

% !TeX root = ../Book4.tex
\begin{figure}[htbp]
	\centering
	\begin{tikzpicture}[scale=.78,>=stealth]
		\begin{scope}
			\draw[->,gray] (-.35,0)--(3.6,0) node[right] {\(x_1\)};
			\draw[->,gray] (0,-.35)--(0,3.55) node[above] {\(x_2\)};
			\filldraw[fill=blue!10,draw=blue!65!black,thick] (0,0)--(3,0)--(0,3)--cycle;
			\node[below left] at (0,0) {\(0\)};
			\node[below] at (3,0) {\(s\)};
			\node[left] at (0,3) {\(s\)};
			\node at (1.45,-.8) {单纯形};
		\end{scope}
		\begin{scope}[xshift=5.2cm]
			\draw[->,gray] (-.35,0)--(3.6,0) node[right] {\(x_1\)};
			\draw[->,gray] (0,-.35)--(0,3.55) node[above] {\(x_2\)};
			\filldraw[fill=green!10,draw=green!45!black,thick] (0,0) rectangle (3,2.3);
			\node[below left] at (0,0) {\(0\)};
			\node[below] at (3,0) {\(a\)};
			\node[left] at (0,2.3) {\(b\)};
			\node at (1.45,-.8) {区间乘积};
		\end{scope}
		\begin{scope}[xshift=10.4cm]
			\draw[->,gray] (-.35,0)--(3.6,0) node[right] {\(x_1\)};
			\draw[->,gray] (0,-.35)--(0,3.55) node[above] {\(x_2\)};
			\filldraw[fill=orange!12,draw=orange!70!black,thick] (1,0)--(3,0)--(0,3)--(0,1)--cycle;
			\draw[dashed,gray] (1,0)--(0,0)--(0,1);
			\draw[very thick,red!70!black] (1,0)--(0,1);
			\node[below] at (1,0) {\(\varepsilon\)};
			\node[left] at (0,1) {\(\varepsilon\)};
			\node at (1.45,-.8) {在顶点截角};
		\end{scope}
	\end{tikzpicture}
	\caption{三类 Delzant 多边形。截角将原顶点替换成一条新边；新边对应吹起中的例外除子。}
	\label{g4:fig:1203:polytopes}
\end{figure}


\subsection{从多面体构造辛约化}
\label{g4:subsec:1203:02}

给定式\eqref{g4:eq:1203:facet-presentation}中的 Delzant 多面体，定义整数线性映射
\begin{equation}\label{g4:eq:1203:normal-map}
	\beta:\mathbb R^d\longrightarrow\mathfrak t,\qquad
	\beta(e_j)=\nu_j.
\end{equation}
在任一顶点挑出的 \(n\) 个 \(\nu_j\) 已经构成整数基，所以
\(\beta(\mathbb Z^d)=N\)。它诱导满射环面同态
\[
	\overline\beta:\mathbb T^d\longrightarrow T,\qquad
	K=\ker\overline\beta,\qquad \mathfrak k=\ker\beta.
\]

\begin{lemma}\label{g4:lem:1203:kernel-torus}
	\(K\) 是连通的 \((d-n)\) 维环面，且有正合列
	\[
		0\longrightarrow\mathfrak k
		\xrightarrow{\ i\ }\mathbb R^d
		\xrightarrow{\ \beta\ }\mathfrak t\longrightarrow0,
		\qquad
		0\longrightarrow\mathfrak t^*
		\xrightarrow{\ \beta^*\ }(\mathbb R^d)^*
		\xrightarrow{\ i^*\ }\mathfrak k^*\longrightarrow0.
	\]
\end{lemma}
\begin{proof}
	选顶点 \(v\)。把 \(N\) 中的向量按 \(\{\nu_j:j\in I_v\}\) 展开，并把系数放回 \(\mathbb Z^d\) 的对应位置，得到整数右逆 \(a:N\to\mathbb Z^d\)，满足 \(\beta a=I\)。于是
	\[
		\mathbb Z^d=(\ker\beta\cap\mathbb Z^d)\oplus a(N),
		\qquad
		\mathbb R^d=\mathfrak k\oplus a(\mathfrak t).
	\]
	取 \(2\pi\) 倍整数格的商后，
	\[
		\mathbb T^d\cong
		\bigl(\mathfrak k/(2\pi(\mathfrak k\cap\mathbb Z^d))\bigr)\times T.
	\]
	在这个分解中 \(\overline\beta\) 是第二因子投影，所以其核恰为第一个连通环面。两条线性正合列由核、像及线性对偶得到。
\end{proof}

在 \(\mathbb C^d\) 上取 \(\omega_0=\sum_j\mathrm{d}x_j\wedge\mathrm{d}y_j\)，让 \(\mathbb T^d\) 按正向逐坐标旋转
\[
	t\cdot z=(e^{it_1}z_1,\ldots,e^{it_d}z_d)
\]
作用。其矩映射与平移后的矩映射为
\[
	\mu_{\mathbb C}(z)=-\frac12(|z_1|^2,\ldots,|z_d|^2),\qquad
	\phi(z)=\mu_{\mathbb C}(z)+\lambda,\quad
	\lambda=(\lambda_1,\ldots,\lambda_d).
\]
正旋转的基本场为 \(-y_j\partial_{x_j}+x_j\partial_{y_j}\)，插入标准辛形式得到
\(-\mathrm{d}(|z_j|^2/2)\)，所以这里的负号与本章约定一致。限制到 \(K\)，矩映射为 \(i^*\phi\)。定义
\begin{equation}\label{g4:eq:1203:delzant-level}
	Z_P=(i^*\phi)^{-1}(0),\qquad M_P=Z_P/K.
\end{equation}
先把它们看作水平集和拓扑商；光滑性将在下一小节证明。

\begin{lemma}\label{g4:lem:1203:level-moduli}
	对每个 \(z\in Z_P\)，存在唯一 \(x\in P\) 满足
	\begin{equation}\label{g4:eq:1203:moduli-equations}
		\beta^*x=\phi(z),\qquad
		\frac{|z_j|^2}{2}=\lambda_j-\langle x,\nu_j\rangle.
	\end{equation}
	反过来，每个 \(x\in P\) 都可由这些模长及任意相位得到 \(Z_P\) 中的点。特别地，\(\phi(Z_P)=\beta^*(P)\)。
\end{lemma}
\begin{proof}
	\(i^*\phi(z)=0\) 与对偶正合列说明
	\(\phi(z)\in\operatorname{im}\beta^*\)，单射性给出唯一 \(x\)。
	逐分量写出 \(\beta^*x=\phi(z)\)，就是第二组等式。其右端非负，恰好说明 \(x\) 满足 \(P\) 的全部不等式。反之，对给定 \(x\in P\)，取
	\(z_j=\sqrt{2(\lambda_j-\langle x,\nu_j\rangle)}\)
	并任意附加相位，就满足所需方程。
\end{proof}

\subsection{自由性和紧性的核验}
\label{g4:subsec:1203:03}

一个坐标 \(z_j\) 为零，当且仅当 \(x\) 落在第 \(j\) 个面上。因此模长方程不仅恢复多面体，也把可能的稳定子与相交的面联系起来。

\begin{proposition}\label{g4:prop:1203:free-compact-level}
	\(Z_P\) 非空、紧且连通，\(K\) 在 \(Z_P\) 上自由。零是 \(i^*\phi\) 的正则值，故 \(Z_P\) 为实维数 \(d+n\) 的光滑无边界子流形。
\end{proposition}
\begin{proof}
	非空性由\cref{g4:lem:1203:level-moduli}得到。每个函数
	\(h_j(x)=\lambda_j-\langle x,\nu_j\rangle\) 在紧多面体 \(P\) 上有界，所以模长方程给出所有 \(|z_j|\) 的共同上界。水平集本身闭，因此紧。映射
	\[
		P\times\mathbb T^d\longrightarrow Z_P,\qquad
		(x,t)\longmapsto
		\bigl(\sqrt{2h_j(x)}\,e^{it_j}\bigr)_{j=1}^d
	\]
	连续且满射。定义域连通，故 \(Z_P\) 连通；面上的平方根只需在这里连续。

	固定 \(z\in Z_P\)，令 \(x\) 为其模长对应点，并记
	\(J(x)=\{j:\langle x,\nu_j\rangle=\lambda_j\}\)。
	逐坐标旋转的稳定子是只在 \(J(x)\) 中允许非平凡相位的子环面 \(\mathbb T^{J(x)}\)。取包含 \(x\) 的最小面上的一个顶点 \(v\)，则 \(J(x)\subset I_v\)。由顶点的整数基条件，
	\(\overline\beta:\mathbb T^{I_v}\to T\)
	为同构。因此
	\[
		K_z=K\cap\mathbb T^{J(x)}
		\subseteq K\cap\mathbb T^{I_v}=\{1\}.
	\]
	这证明作用自由，包含有限稳定子在内都已被排除。

	由\cref{g4:prop:1002:moment-differential}，自由性使 \(K\) 的无穷小稳定子为零，故 \(\mathrm{d}(i^*\phi)\) 在 \(Z_P\) 上满射。零因而正则，水平集维数为
	\(2d-(d-n)=d+n\)。
\end{proof}

\(K\) 紧，作用自动固有。于是\cref{g4:thm:1002:symplectic-reduction}赋予 \(M_P=Z_P/K\) 唯一辛形式 \(\omega_P\)，满足
\begin{equation}\label{g4:eq:1203:reduced-form}
	\pi^*\omega_P=\jmath^*\omega_0,\qquad
	\pi:Z_P\to M_P,\quad \jmath:Z_P\hookrightarrow\mathbb C^d.
\end{equation}
它是紧连通的 \(2n\) 维无边界辛流形。构造中，简单性和光滑格条件保证自由性；多面体的有界性保证紧性；正则性则由自由性与矩映射微分判据得到。

\subsection{矩映射像的恢复}
\label{g4:subsec:1203:04}

逐坐标旋转保持 \(Z_P\)，并与 \(K\) 交换，因而在商上留下
\(\mathbb T^d/K\cong T\) 的作用。模长方程中的 \(x\) 沿 \(K\) 轨道不变，所以给出光滑映射
\begin{equation}\label{g4:eq:1203:residual-moment}
	\mu_P:M_P\longrightarrow\mathfrak t^*,\qquad
	\beta^*(\mu_P\pi)=\phi|_{Z_P}.
\end{equation}
其光滑性可由线性单射 \(\beta^*\) 在其像上的逆和商映射的局部截面看出。

\begin{theorem}[Delzant 构造]\label{g4:thm:1203:delzant-construction}
	设 \(P\subset\mathfrak t^*\) 是关于格 \(N\) 的 \(n\) 维 Delzant 多面体。上述构造给出紧连通辛流形 \((M_P^{2n},\omega_P)\) 及 \(T=\mathfrak t/(2\pi N)\) 的有效 Hamilton 作用。\(\mu_P\) 是等变矩映射，且
	\[
		\mu_P(M_P)=P.
	\]
\end{theorem}
\begin{proof}
	只需验证剩余作用的矩映射方程和有效性。给定 \(\xi\in\mathfrak t\)，选 \(\widetilde\xi\in\mathbb R^d\) 满足 \(\beta\widetilde\xi=\xi\)。则
	\[
		\pi^*\langle\mu_P,\xi\rangle
		=\langle\phi,\widetilde\xi\rangle|_{Z_P}.
	\]
	利用\eqref{g4:eq:1203:reduced-form}和相关基本场，对两边求微分得到
	\[
		\pi^*\mathrm{d}\langle\mu_P,\xi\rangle
		=\jmath^*(\iota_{\widetilde\xi_{\mathbb C^d}}\omega_0)
		=\pi^*(\iota_{\xi_{M_P}}\omega_P).
	\]
	改变提升只增加 \(K\) 的方向，该方向在限制形式的核中；所以这个计算与提升选择无关。淹没拉回单射，得到矩映射方程。\(\mu_P\) 由模长给出，故对 \(T\) 不变，也就是环面情形的等变性。像等式已由\cref{g4:lem:1203:level-moduli}证明。

	取 \(x\in\operatorname{int}P\)，相应的 \(z\in Z_P\) 全部坐标非零。若 \(t\in T\) 固定 \([z]\)，取其提升 \(g\in\mathbb T^d\)，存在 \(k\in K\) 使 \(gz=kz\)。因逐坐标旋转在这个 \(z\) 处的稳定子平凡，有 \(g=k\)，故 \(t=1\)。于是作用有自由轨道，特别有效。
\end{proof}

\begin{proposition}\label{g4:prop:1203:orbit-fibres}
	对任意 \(x\in P\)，\(\mu_P^{-1}(x)\) 恰为一条 \(T\) 轨道，其稳定子为
	\[
		T_x=\overline\beta(\mathbb T^{J(x)}).
	\]
	若 \(x\) 在一个余维 \(r\) 的面的相对内部，则 \(T_x\) 为 \(r\) 维子环面。特别地，顶点对应唯一固定点，内部对应自由轨道，且 \(\mu_P\) 诱导同胚 \(M_P/T\cong P\)。
\end{proposition}
\begin{proof}
	固定 \(x\) 后各坐标模长全部固定，所有点只差相位，故 \(\mathbb T^d\) 在其 \(Z_P\) 原像上传递。取 \(K\) 商后，剩余 \(T\) 仍传递。某个 \(g\in\mathbb T^d\) 固定 \([z]\)，等价于 \(g\in K\mathbb T^{J(x)}\)，所以剩余稳定子正是所列像。由前面的基底包含关系，\(\overline\beta\) 在 \(\mathbb T^{J(x)}\) 上单射。简单性给出 \(|J(x)|=r\)，故该稳定子为 \(r\) 维子环面。最后，连续诱导映射 \(M_P/T\to P\) 是从紧空间到 Hausdorff 空间的双射，因而为同胚。
\end{proof}

\begin{proposition}\label{g4:prop:1203:vertex-chart}
	固定顶点 \(v\)，把其面指标重排成 \(I_v=\{1,\ldots,n\}\)，并令 \(m_i\) 为外法向量 \(\nu_i\) 的对偶基。置
	\[
		x(w)=v-\sum_{i=1}^n\frac{|w_i|^2}{2}m_i,\qquad
		D_v=\{w\in\mathbb C^n:h_j(x(w))>0\text{ 对所有 }j\notin I_v\}.
	\]
	则
	\[
		w\longmapsto
		\left[
		w_1,\ldots,w_n,\sqrt{2h_{n+1}(x(w))},\ldots,\sqrt{2h_d(x(w))}
		\right]
	\]
	是从 \(D_v\) 到 \(M_P\) 中相应开集的辛坐标参数化，且矩映射为 \(x(w)\)。特别地，顶点固定点附近的形式严格为
	\(\sum_i\mathrm{d}(\operatorname{Re}w_i)\wedge\mathrm{d}(\operatorname{Im}w_i)\)。
\end{proposition}
\begin{proof}
	顶点基条件使 \(K\to\mathbb T^{I_v^c}\) 的坐标投影为同构：任给非活动位置的相位，用 \(\beta|_{\mathbb Z^{I_v}}\) 的整数逆唯一补出活动位置的相位，使整体属于 \(K\)。因此当全部非活动坐标非零时，可以唯一用 \(K\) 将它们转为正实数。剩余活动坐标为任意 \(w_i\)，模长方程对活动面给出 \(x=x(w)\)，再对非活动面给出所列正平方根。这证明双向光滑且彼此为逆。

	在这个截面上，每个非活动复坐标的虚部恒为零，故它对 \(\omega_0\) 的贡献为零；活动坐标直接贡献标准形式。约化拉回等式于是给出所述辛形式和矩映射。
\end{proof}

例如 \(P=[0,a]\)、\(a>0\)，可取外法向量 \((-1,1)\) 与支撑常数 \((0,a)\)。此时 \(K\) 为 \(\mathbb T^2\) 中的对角圆，水平方程为
\[
	|z_1|^2+|z_2|^2=2a.
\]
约化空间是 \(\mathbb{CP}^1\)，由\cref{g4:thm:1003:projective-reduction}，其辛形式为 \(2a\,\omega_{\mathrm{FS}}\)，总面积为 \(2\pi a\)。剩余矩坐标为 \(x=|z_1|^2/2\)，取值恰为 \([0,a]\)。这个例子也固定了本章中格长度与辛面积之间的 \(2\pi\) 比例。

\ProseParagraphBreak

以上完成了从多面体到辛流形的存在性构造。一般构造与区间例子可对照 Cannas da Silva 的第6.5节。\cite[第174--178页]{Cannas2006Symplectic}要从任意有效半维 Hamilton 环面作用反推出同一个模型，还需要比较局部标准形并粘合；这是下一节的任务。
% 纲要条目排为编号小节标题，并进入目录与书签；说明段仍保留于纲要注释中。

% !TeX root = ../../Book4.tex
% 纲要标题：### 12.4 分类的唯一性
% 纲要位置：几何学计划纲要.md，第 3947 行。
% 正文状态：初稿；尚未正式审核。

% 纲要内容开始
% 
% * 局部标准模型的比较
% * 自由轨道部分的粘合
% * 边界轨道的延拓
% * 等变辛同构与平移自由
% 
% 纲要内容结束

\section{分类的唯一性}
\label{g4:sec:1204}

由多面体作约化得到一个流形，还没有说明给定的流形一定与它相同。唯一性需要比较两套局部模型，并把局部同构粘成整体同构。多面体内部的纤维是自由环面轨道；边界上有圆方向收缩，角坐标会失效。因此我们先确定全部稳定子和边界模型，再把粘合误差写成底空间上的光滑 Hamilton 函数。凸多面体的可缩性将消去这些误差。

\subsection{局部标准模型的比较}
\label{g4:subsec:1204:01}
本节固定 \(T=\mathfrak t/(2\pi\Lambda)\)，其中 \(\Lambda\cong\mathbb Z^n\)。辛环面流形是带有效 Hamilton \(T\) 作用的紧连通辛流形 \((M^{2n},\omega)\)，并指定矩映射 \(\mu:M\to\mathfrak t^*\)。先由一般轨道标准形推出半维作用的特殊性质。

\begin{proposition}\label{g4:prop:1204:toric-orbits}
	设 \((M^{2n},\omega,T,\mu)\) 为辛环面流形。则每个稳定子 \(H=T_x\) 都连通。若 \(\dim H=h\)，辛切片的复维数为 \(h\)，其 \(h\) 个权构成 \(H\) 的特征格的一组整数基。每条矩纤维恰是一条 \(T\) 轨道，矩映射诱导同胚
	\[
		M/T\xrightarrow{\ \cong\ }P=\mu(M).
	\]
	\(P\) 是满维 Delzant 多面体；余维 \(h\) 面的相对内部对应稳定子维数为 \(h\) 的轨道。
\end{proposition}
\begin{proof}
	\Cref{g4:thm:1201:orbit-normal-form}给出 \(x\) 的轨道附近的模型
	\[
		T\times_H(\mathfrak h^\circ\times V),\qquad
		\dim_{\mathbb R}V=2n-2\dim(T/H)=2h.
	\]
	给 \(V\) 取 \(H\) 不变的相容复结构。紧交换群的酉表示可同时对角化，因而有 \(h\) 个字符 \(\chi_j:H\to S^1\)。它们给出的同态 \(\chi:H\to(S^1)^h\) 是单射。事实上，核中的元素在 \(V\) 上恒等，因 \(T\) 交换，它在模型的轨道方向及 \(\mathfrak h^\circ\) 上也恒等，所以在 \(M\) 的一个开集上恒等。

	为说明这迫使该元素全局恒等，平均任一 Riemann 度量，使 \(T\) 作用等距。第三卷定理2.1.2的测地初值唯一性使等距映射与指数映射交换，命题2.1.4又说明指数映射在零点附近是局部微分同胚。因此一个等距映射若固定一点且微分在该点为恒等，就在该点的一个邻域恒等。具有这两项性质的点集既闭又开；从已经固定的开集出发，连通性使它等于 \(M\)。作用有效，故该元素只能是单位元。

	于是 \(\chi\) 的微分单射；定义域和值域的 Lie 代数维数都是 \(h\)，故它的像含单位元邻域。紧性使像闭，而 \((S^1)^h\) 连通，像必为整个环面。因此 \(H\cong(S^1)^h\)，特别地 \(H\) 连通，且各 \(\chi_j\) 构成特征格的整数基。当 \(h=0\) 时，同一论证说明稳定子平凡。

	选择 \(H\) 的整数格补，局部模型可写为
	\begin{equation}\label{g4:eq:1204:corner-model}
		T^{n-h}\times\mathbb R^{n-h}_p\times\mathbb C^h_z,
		\qquad
		\omega=\sum_{a=1}^{n-h}\mathrm{d}\theta_a\wedge\mathrm{d}p_a
		+\sum_{j=1}^h\mathrm{d}x_j\wedge\mathrm{d}y_j,
	\end{equation}
	其中 \(H\) 以各坐标的正旋转作用，\(z_j=x_j+\mathrm{i}y_j\)。相对于相应的整数基，矩变量为 \((p,-|z_1|^2/2,\ldots,-|z_h|^2/2)\)，另加 \(\mu(x)\) 的平移。一个矩值确定全部 \(p_a\) 和 \(|z_j|\)，因而局部矩纤维恰为一条轨道。这也给出局部轨道商的象限坐标。

	每条轨道因此在其矩纤维中相对开。\Cref{g4:thm:1202:connected-fibres}说明矩纤维连通，故不能分成两条以上互不相交的开轨道，纤维只能是一条轨道。\(T\) 紧，轨道商是紧 Hausdorff 空间；诱导的连续双射到 \(P\subset\mathfrak t^*\) 因而为同胚。

	\Cref{g4:thm:1202:convexity}已给出 \(P\) 为凸多面体。局部模型中的 \(h\) 个半轴与 \(n-h\) 个自由坐标说明它满维，并在该点有恰好 \(h\) 个相交的边界面。顶点恰好对应 \(h=n\)。此时向内部的原始边向量是切片权的相反数，构成整数基；等价地，各面的原始外法向也构成对偶整数基。这正是 Delzant 条件。一般面上的稳定子由对应边界面的法向圆群生成。
\end{proof}

两个辛环面流形若有同一个矩多面体 \(P\)，上面的结论在每个 \(b\in P\) 给出相同的稳定子、切片权及局部矩锥。一般轨道标准形遂给出 \(b\) 上方轨道附近保持矩映射的 \(T\) 等变辛同构。我们可以把它的定义域取为 \(\mu^{-1}(U)\)，其中 \(U\) 是 \(P\) 中的相对开邻域：若一个轨道邻域包含 \(\mu^{-1}(b)\)，其补集紧，补集的矩像不含 \(b\)，所以缩小 \(U\) 后整条邻近纤维都落在该轨道邻域中。以后取有限个这样的小凸邻域覆盖 \(P\)。

\ProseParagraphBreak

这里的局部同构保存固定的环面 \(T\)。在顶点改用整数格基只是表示同一个局部模型，不能把不同顶点处任意的实线性换基当成全局允许的群变换。

\subsection{自由轨道部分的粘合}
\label{g4:subsec:1204:02}
先在 \(P\) 的内部观察粘合误差。此处作用自由，\(\mu\) 的分量是全局给定的作用变量。在一个小凸集 \(U\subset\operatorname{int}P\) 上，第十一章给出
\[
	\mu^{-1}(U)\cong T\times U,\qquad
	\omega=\sum_{j=1}^n\mathrm{d}\theta_j\wedge\mathrm{d}b_j,
\]
其中 \(b=(b_1,\ldots,b_n)\) 是固定的矩坐标。与 \(T\) 作用交换且保持 \(b\) 的映射只能是
\begin{equation}\label{g4:eq:1204:phase-translation}
	(\theta,b)\longmapsto(\theta+a(b),b),\qquad a:U\to T.
\end{equation}
在 \(U\) 上取 \(a\) 的实提升 \(\widetilde a:U\to\mathfrak t\)。代入辛形式，映射保辛当且仅当
\begin{equation}\label{g4:eq:1204:closed-phase}
	\sum_j\mathrm{d}\widetilde a_j\wedge\mathrm{d}b_j=0,
	\quad\text{即}\quad
	\alpha=\sum_j\widetilde a_j\,\mathrm{d}b_j\text{ 闭}.
\end{equation}
凸性给出 \(\alpha=\mathrm{d}f\)，因此该映射是 \(f\circ\mu\) 的时间 \(1\) Hamilton 流。

\ProseParagraphBreak

这个表达把局部同构之间的误差变成了可相加的函数。记
\begin{equation}\label{g4:eq:1204:hamilton-automorphism}
	\Phi_f(x)=\exp\bigl(\mathrm{d}f_{\mu(x)}\bigr)\cdot x,
\end{equation}
这里用 \(\mathrm{T}_b^*\mathfrak t^*\cong\mathfrak t\) 识别微分，\(\exp:\mathfrak t\to T\) 的核为 \(2\pi\Lambda\)。有 \(\Phi_f\Phi_g=\Phi_{f+g}\)。在连通 \(U\) 上，\(\Phi_f\) 为恒等，当且仅当
\begin{equation}\label{g4:eq:1204:affine-kernel}
	f(b)=c+2\pi\langle b,m\rangle,
	\qquad c\in\mathbb R,\quad m\in\Lambda.
\end{equation}
必要性来自自由性：\(\mathrm{d}f\) 只能取在离散集 \(2\pi\Lambda\) 中，连续性使它恒定；充分性则由指数映射直接得到。

若只在 \(\operatorname{int}P\) 处理这些函数，所得同构还未必在收缩的角圆处光滑。下一小节将证明，光滑的边界模型恰好保证相位及其生成函数也能光滑延至多面体各面。

\subsection{边界轨道的延拓}
\label{g4:subsec:1204:03}
设 \(U\) 为多面体的相对开集，\(f\) 为 \(U\) 上的函数。若 \(f\) 在每一点附近都是环境 \(\mathfrak t^*\) 上某个光滑函数的限制，则称 \(f\) 在 \(U\) 上光滑，记为 \(f\in C^\infty(U)\)。局部象限中的各阶单侧导数也可描述同一个光滑结构。下面先解释从复坐标到象限坐标所需的光滑性事实。

\begin{lemma}\label{g4:lem:1204:radial-division}
	一个在 \(\mathbb C^h\times\mathbb R^k\) 的原点附近光滑、对各复坐标的独立旋转不变的函数，可局部写成
	\[
		A(|z_1|^2,\ldots,|z_h|^2,p),
	\]
	其中 \(A\) 在象限上光滑。若复值光滑函数 \(G_j\) 只在第 \(j\) 个圆作用下具有权 \(1\)，在其他圆作用下不变，则
	\[
		G_j(z,p)=z_j A_j(|z_1|^2,\ldots,|z_h|^2,p)
	\]
	且 \(A_j\) 同样光滑。
\end{lemma}
\begin{proof}
	先处理一个复变量，并允许其余变量作参数。旋转不变函数在实轴的限制 \(g(r,p)\) 是光滑偶函数，其值决定整个函数。对任意 \(N\)，带参数的 Taylor 公式给出
	\[
		g(r,p)=\sum_{j=0}^{N}a_j(p)r^{2j}+r^{2N+2}R_N(r,p),
	\]
	其中余项及任意预先指定的有限阶参数导数局部有界。令 \(s=r^2\)，并对 \(s>0\) 反复使用 \(\partial_s=(2r)^{-1}\partial_r\)。先选足够高阶的 Taylor 展开，就能验证任意有限阶 \(s,p\) 导数都连续延至 \(s=0\)，其边界值为相应 Taylor 系数。故 \(A(s,p)=g(\sqrt s,p)\) 在闭半轴上各阶光滑。

	这也确实给出环境光滑延拓。可取固定截断函数 \(\chi\)，对给定的边界系数逐项构造
	\[
		\sum_{j\ge0}a_j(p)s^j\chi(s/\varepsilon_j).
	\]
	在一个较小紧参数邻域上，逐次选 \(\varepsilon_j\downarrow0\)，使第 \(j\) 项在所有总阶小于 \(j\) 的导数范数中不超过 \(2^{-j}\)。有限低阶项无此限制。该级数在双侧光滑且具有指定的全部边界导数；从 \(A\) 减去它后得到在 \(s=0\) 各阶为零的函数，将差值在 \(s<0\) 延为零即得光滑延拓。这一构造局部适用于参数。

	对权 \(1\) 的 \(G\)，实轴限制满足 \(G(-r,p)=-G(r,p)\)，且 \(G(0,p)=0\)。公式
	\[
		\frac{G(r,p)}r=\int_0^1\partial_rG(tr,p)\,\mathrm{d}t
	\]
	把商光滑延至零；商为偶函数，由前段写成 \(A(r^2,p)\)。旋转等变性再给出 \(G(z,p)=zA(|z|^2,p)\)。对多个复变量，依次应用同一带参数论证；各次取商和平方径向替换均保持其他旋转的不变性，得到所述结论。
\end{proof}

\begin{proposition}\label{g4:prop:1204:boundary-automorphisms}
	设 \(U\subset P\) 是足够小的凸相对开集。任一保持矩映射的 \(T\) 等变辛自同构 \(a:\mu^{-1}(U)\to\mu^{-1}(U)\) 都可写成 \(a=\Phi_f\)，其中 \(f\in C^\infty(U)\)。两个这样的函数给出同一个映射，当且仅当它们相差式\eqref{g4:eq:1204:affine-kernel}中的函数。
\end{proposition}
\begin{proof}
	在内部自由轨道上，相位 \(g(b)\in T\) 唯一确定。考察式\eqref{g4:eq:1204:corner-model}的一张边界图。等变性和保持矩值说明，映射的自由角部分是 \(\theta\mapsto\theta+u(|z|^2,p)\)，复坐标部分满足
	\[
		z_j\longmapsto z_j A_j(|z|^2,p),\qquad |A_j|=1.
	\]
	这里自由角的圆值相位及 \(A_j\) 的光滑性由\cref{g4:lem:1204:radial-division}得到；\(|A_j|=1\) 先在 \(z_j\ne0\) 由模长不变得到，再以连续性延至 \(z_j=0\)。所以内部的 \(g\) 光滑延至全部边界。两张图得到的延拓在稠密内部相同，故在交集上相同。这证明 \(g:U\to T\) 是全局光滑的。

	凸集上的圆值映射有光滑实提升。具体地，对各分量 \(g_j\)，闭一形式 \(-\mathrm{i}g_j^{-1}\mathrm{d}g_j\) 沿从一个内点到 \(b\) 的线段积分，得到 \(\widetilde g_j(b)\)，再调整其在基点的常数，就有 \(e^{\mathrm{i}\widetilde g_j}=g_j\)。线段积分在各面也光滑。内部的保辛条件为式\eqref{g4:eq:1204:closed-phase}，光滑性使它在边界仍成立。再次沿线段积分，得到 \(\sum_j\widetilde g_j\,\mathrm{d}b_j=\mathrm{d}f\)。故 \(a=\Phi_f\) 在内部相同，并由连续性在全部 \(\mu^{-1}(U)\) 相同。

	反过来，\(f\circ\mu\) 在流形上光滑，其 Hamilton 场沿紧轨道切向，故时间 \(1\) 流存在，并等于式\eqref{g4:eq:1204:hamilton-automorphism}。核的刻画在内部已证，函数的连续性将同一仿射表达延至边界。
\end{proof}

边界处单独一条轨道不能唯一确定相位，因为稳定子中的元素对它作用平凡；邻近自由轨道及光滑延拓却确定了同一相位函数。这个区别使我们能够把全部粘合数据放在多面体上处理。

\subsection{等变辛同构与平移自由}
\label{g4:subsec:1204:04}
还需要一条关于凸集覆盖的具体代数事实。它将消去局部 Hamilton 函数在三重交集上的整数仿射误差。

\begin{lemma}\label{g4:lem:1204:convex-cech}
	设 \(P\) 为非空凸多面体，\(\mathcal U=(U_i)\) 为有限相对开凸覆盖，\(K\) 为任意阿贝尔群。以局部常值 \(K\) 系数构成的这份覆盖的 Čech 复形，在正次数的上同调为零。特别地，每个常值 \(K\) 系数的二上圈都是一上链的余边界。
\end{lemma}
\begin{proof}
	每个非空有限交 \(U_{i_0\cdots i_p}\) 都凸。考虑双复形
	\[
		B^{p,q}=\prod_{i_0<\cdots<i_p}C^q(U_{i_0\cdots i_p};K),
	\]
	横向为交替限制，纵向为奇异上边界，总微分取 \(\delta_{\mathrm C}+(-1)^p\delta_{\mathrm s}\)。对固定交集，沿线段向其中一点作锥，给出增广奇异链的收缩；对偶后可见纵向上同调只在零次为 \(K\)。第二卷引理7.2.2的第一象限比较于是将总上同调识别为所述常值 Čech 上同调。

	另一方面，用覆盖小奇异上链增广各横行。对每条小单形 \(\sigma\)，选一个容纳其全部像的指标 \(j(\sigma)\)。在上链的交替指标首位插入 \(j(\sigma)\)，逐项求值即给出横向收缩；交替和展开后满足 \(\delta h+h\delta=1\)。因此同一比较引理把总上同调识别为小奇异上链的上同调。这里无须令小单形的指标选择随 \(\sigma\) 连续变化。

	第二卷定理2.3.8说明小链包含在整数同调上是同构；小链和全部奇异链都为自由整数链复形，第二卷定理4.2.3的自然上同调万有系数列因而把此同构转成任意 \(K\) 系数的上同调同构。最后，\(P\) 自身也可沿线段收缩到一点，锥算子使其正次上同调为零。结合两种比较，即得结论。这个证明只用了凸性和相对开覆盖，所以也适用于 \(P\) 的边界与角点。
\end{proof}

\begin{theorem}[Delzant 分类的唯一性]\label{g4:thm:1204:uniqueness}
	设 \((M_i^{2n},\omega_i,T,\mu_i)\)、\(i=1,2\)，为两个辛环面流形，环面及其整数格已固定。若 \(\mu_1(M_1)=\mu_2(M_2)=P\)，则存在 \(T\) 等变辛同胚 \(\Psi:M_1\to M_2\)，使 \(\mu_2\Psi=\mu_1\)。
\end{theorem}
\begin{proof}
	由\cref{g4:prop:1204:toric-orbits}及其后的局部比较，取有限相对开凸覆盖 \((U_i)\) 和局部同构
	\[
		\psi_i:\mu_1^{-1}(U_i)\longrightarrow\mu_2^{-1}(U_i).
	\]
	这些集合可取为 \(P\) 与小 Euclid 球的交，因而全部非空有限交仍凸。对两重交集，定义源流形上的转移自同构 \(a_{ij}=\psi_i^{-1}\psi_j\)。\Cref{g4:prop:1204:boundary-automorphisms}在局部给出光滑相位；其在内部的唯一性使这些相位在整个交集相容。交集凸，故同一命题的积分步骤在整个交集给出 \(f_{ij}\in C^\infty(U_i\cap U_j)\)，满足
	\[
		a_{ij}=\Phi_{f_{ij}},\qquad f_{ji}=-f_{ij}.
	\]
	转移映射满足 \(a_{ij}a_{jk}=a_{ik}\)，所有这类自同构交换，因此
	\[
		f_{ij}+f_{jk}-f_{ik}
		=c_{ijk}+2\pi\langle b,m_{ijk}\rangle,
		\qquad (c_{ijk},m_{ijk})\in\mathbb R\oplus\Lambda.
	\]
	每个交集连通，右端的系数为常量。它们是一个 \(K=\mathbb R\oplus\Lambda\) 系数的二上圈，因为左端是 \(f\) 的 Čech 余边界，而 \(\delta^2f=0\)。仿射函数在非空交集中的表达唯一：每个这样的相对开交集都含有 \(P\) 的内点，不能使一个非零仿射函数恒为零。

	由\cref{g4:lem:1204:convex-cech}，可找到常整数仿射的一上链，其余边界恰等于上述误差。从 \(f_{ij}\) 减去该一上链不会改变 \(\Phi_{f_{ij}}\)，却使
	\begin{equation}\label{g4:eq:1204:exact-function-cocycle}
		f_{ij}+f_{jk}=f_{ik}.
	\end{equation}
	这一步同时处理实常数和整数周期，因而后续可以对真正的函数余圈使用单位分解。

	在环境 Euclid 空间取从属于这些小邻域的光滑非负截断，并在 \(P\) 附近以它们的正和归一化，得到 \(P\) 上的光滑单位分解 \((\rho_k)\)，且 \(\operatorname{supp}\rho_k\subset U_k\)。令
	\[
		h_i=\sum_k\rho_k f_{ik}\quad\text{于 }U_i.
	\]
	每项在相应交集定义，再借助支集条件补零；有限求和保证光滑。式\eqref{g4:eq:1204:exact-function-cocycle}给出
	\[
		h_i-h_j=\sum_k\rho_k(f_{ik}-f_{jk})=f_{ij}.
	\]
	所以修正后的局部同构 \(\widetilde\psi_i=\psi_i\Phi_{h_i}\) 满足
	\[
		\widetilde\psi_i^{-1}\widetilde\psi_j
		=\Phi_{-h_i}\Phi_{f_{ij}}\Phi_{h_j}=\operatorname{id}.
	\]
	它们拼成全局 \(\Psi\)。每张图中的映射光滑、保辛、等变并保持矩映射，这些性质在拼接后保持。每条目标纤维上，它由任一覆盖该底点的局部同构给出双射，所以 \(\Psi\) 全局双射；局部逆也相容，故为所述等变辛同胚。
\end{proof}

结合前节的构造，得到整个分类：固定 \(T\)、指定矩映射的紧连通辛环面流形，按保持矩映射的 \(T\) 等变辛同胚分类，与 \(\mathfrak t^*\) 中的 Delzant 多面体一一对应。分类定理的标准陈述可参见 Cannas da Silva 的定理6.8；前节的约化与本节的粘合分别完成存在性和唯一性。\cite[第6.4节，第169--174页]{Cannas2006Symplectic}

\begin{proposition}\label{g4:prop:1204:translation-lattice-freedom}
	在上述分类中，若不指定矩映射的加法常数，则多面体按平移识别。若还允许以环面自同构重标作用，则多面体进一步按保持整数格的线性变换识别，即按整数仿射变换分类。
\end{proposition}
\begin{proof}
	同一个连通辛流形上，同一作用的两个矩映射对每个 \(\xi\in\mathfrak t\) 的分量有相同微分。因此它们的差是常向量 \(c\in\mathfrak t^*\)，矩像相差平移。

	若环面自同构 \(\kappa:T\to T\) 的微分为 \(A\)，则 \(A\Lambda=\Lambda\)，相对于整数基有 \(A\in\operatorname{GL}(n,\mathbb Z)\)。若 \(\Psi(t\cdot x)=\kappa(t)\cdot\Psi(x)\) 且 \(\Psi\) 保辛，则
	\[
		\mathrm{d}\langle\mu_2\Psi,A\xi\rangle
		=\Psi^*\iota_{(A\xi)_{M_2}}\omega_2
		=\iota_{\xi_{M_1}}\omega_1
		=\mathrm{d}\langle\mu_1,\xi\rangle.
	\]
	所以 \(A^{\mathsf T}\mu_2\Psi=\mu_1+c\)。这给出对偶空间中的整数线性变换及平移；反过来，重标群作用及平移矩映射后应用唯一性定理，即得相应同构。
\end{proof}

多面体的尺度并未被上述等价关系消去。例如一维区间长度为 \(\ell\) 时，对应球面的辛面积为 \(2\pi\ell\)，不同长度不能由辛同胚识别。分类保存的是格、各面位置及其相对尺度；只保留多面体的组合类型，会丢失这些辛几何数据。
% 纲要条目排为编号小节标题，并进入目录与书签；说明段仍保留于纲要注释中。

% !TeX root = ../../Book4.tex
% 纲要标题：### 12.5 与代数环面簇的比较
% 纲要位置：几何学计划纲要.md，第 3954 行。
% 正文状态：初稿；尚未正式审核。

% 纲要内容开始
% 
% * 法扇与不变除子
% * 极化和整多面体
% * 射影空间、乘积及吹起模型
% * 一般非光滑多面体的新现象
% 
% ---
% 
% 纲要内容结束

\section{与代数环面簇的比较}
\label{g4:sec:1205}

一个多面体同时包含两种信息：面的法向量规定坐标图怎样粘合，面的相对位置规定辛形式的大小。把圆相位扩展为非零复数的乘法，Delzant 构造便出现代数几何中的单项式图。本节从这些图出发建立比较，只使用仿射与射影坐标、线丛的过渡函数和明确的约化计算。

\subsection{法扇与不变除子}
\label{g4:subsec:1205:01}

沿用\secref{g4:sec:1203}的格 \(N\)、对偶格 \(M_{\mathbb Z}\) 及原始外法向量 \(\nu_j\)。\term[daishuhuanmian]{代数环面}[algebraic torus]取为
\[
	T_{\mathbb C}=\operatorname{Hom}(M_{\mathbb Z},\mathbb C^*)
	\cong(\mathbb C^*)^n.
\]
对 \(m\in M_{\mathbb Z}\)，相应字符记为 \(\chi^m\)；在选定格基后，它是 Laurent 单项式，且 \(\chi^{m+m'}=\chi^m\chi^{m'}\)。单位模长部分就是 \(T\)。

\begin{definition}\label{g4:def:1205:normal-fan}
	对多面体 \(P\) 的非空面 \(F\)，定义\term[fazhui]{法锥}[normal cone]
	\[
		\sigma_F=
		\{u\in\mathfrak t:\langle x,u\rangle
		\text{ 在全部 }x\in F\text{ 处达到其在 }P\text{ 上的最大值}\}.
	\]
	它由包含 \(F\) 的各面的外法向量生成。所有这些锥组成 \(P\) 的\term[fashan]{法扇}[normal fan]，记为 \(\Sigma_P\)。
\end{definition}

这些锥覆盖 \(\mathfrak t\)，因为每个线性函数都在某个面上取最大值。两个法锥的交，是包含相应两面的最小面之法锥；它同时是两个锥的面。可从最大值条件看出交的等式：在线性函数的最大值面中同时包含两个面，就包含它们生成的最小面。反向包含直接成立。面包含关系在法锥处反向；在一个顶点附近用各活动面的线性不等式，就能把这一关系写成删去若干非负系数生成元。因此它们确实按共同的面相接。Delzant 条件说明，每个顶点的最大法锥由 \(N\) 的一组基底生成。

\ProseParagraphBreak

固定顶点 \(v\)，把其外法向基记为 \(\nu_1,\ldots,\nu_n\)，对偶基记为 \(m_1,\ldots,m_n\)。对应的仿射图取为
\[
	U_v=\mathbb C^n,\qquad u_i=\chi^{m_i}.
\]
在稠密环面上，每个字符 \(\chi^m\) 都可写成 \(\prod_i u_i^{\langle m,\nu_i\rangle}\)。它在整张图中正则，当且仅当所有指数非负。用半群记号表示，其坐标环为
\[
	\mathbb C[\sigma_v^\vee\cap M_{\mathbb Z}]
	=\mathbb C[u_1,\ldots,u_n],
	\qquad
	\sigma_v^\vee=\{m:\langle m,u\rangle\ge0\ (u\in\sigma_v)\}.
\]
这里半群环表示以 \(\chi^m\) 为基、按指数相加作乘法的代数。两张顶点图在共同面的对应开集上，通过这些字符的 Laurent 单项式表达粘合；负指数只出现在已经要求非零的坐标上。指数相加保证三重重叠上的变换相容。这是第五卷命题1.3.2的仿射图方法在单项式坐标中的运用。

在一个复流形上，若实闭二形式 \(\omega\) 与其可积复结构 \(J\) 相容，即
\(\omega(Ju,Jv)=\omega(u,v)\) 且 \(\omega(u,Ju)>0\) 对非零 \(u\) 成立，就称它为\term[Kahlerxingshi]{Kähler 形式}[Kähler form]。下面的商图将同时给出复结构与这种形式。

\begin{proposition}\label{g4:prop:1205:complex-quotient}
	设 \(P\) 为 Delzant 多面体，并令
	\[
		\mathcal U=\bigcup_v\mathcal U_v\subset\mathbb C^d,\qquad
		\mathcal U_v=\{z:z_j\ne0\text{ 对所有 }j\notin I_v\}.
	\]
	整数映射 \(\beta\) 诱导复环面满射
	\((\mathbb C^*)^d\to T_{\mathbb C}\)，核记为 \(K_{\mathbb C}\)。
	商 \(\mathcal U/K_{\mathbb C}\) 由上述仿射图 \(U_v\) 粘成，且有自然微分同胚
	\[
		M_P=Z_P/K\longrightarrow\mathcal U/K_{\mathbb C}.
	\]
	在这个复结构下，\(\omega_P\) 是 Kähler 形式。因而辛构造与上述单项式图给出同一个光滑复流形。
\end{proposition}
\begin{proof}
	对 \(z\in Z_P\)，其零坐标集包含于某个 \(I_v\)，所以 \(Z_P\subset\mathcal U\)。顶点整数基使 \(K_{\mathbb C}\to(\mathbb C^*)^{I_v^c}\) 的坐标投影为同构。具体地，若
	\(\nu_j=\sum_{i\in I_v}c_{ij}\nu_i\)，则核条件为
	\[
		k_i=\prod_{j\notin I_v}k_j^{-c_{ij}}\quad(i\in I_v).
	\]
	所以可唯一用 \(K_{\mathbb C}\) 把 \(\mathcal U_v\) 中的非活动坐标全部变成 \(1\)。剩余坐标为
	\[
		u_i=z_i\prod_{j\notin I_v}z_j^{c_{ij}},\qquad i\in I_v.
	\]
	它们可任取复数，给出 \(\mathcal U_v/K_{\mathbb C}\cong\mathbb C^n\)。在剩余复环面作用下，\(u_i\) 的字符正是 \(m_i\)。不同图的重叠只要求某些活动坐标非零，过渡公式便是前述 Laurent 单项式。

	接着证明每个复轨道与 \(Z_P\) 相交于恰好一个 \(K\) 轨道。由整数分裂，
	\(K_{\mathbb C}=K\exp(\mathfrak k)\)，其中 \(\exp(\tau)\) 表示正实数逐坐标伸缩 \(e^{\tau_j}\)。固定 \(z\in\mathcal U_v\)，其零坐标集 \(J\) 包含于 \(I_v\)。在由这些面相交得到的面内取相对内点 \(x\)，使
	\[
		a_j=\lambda_j-\langle x,\nu_j\rangle
		\quad\text{满足}\quad
		a_j=0\ (j\in J),\qquad a_j>0\ (j\notin J).
	\]
	这样的点可从顶点沿不属于 \(J\) 的内向边作足够小的正向移动得到。对 \(\tau\in\mathfrak k\) 考察
	\[
		\mathcal F_z(\tau)
		=\frac14\sum_j|z_j|^2e^{2\tau_j}
			-\sum_j\lambda_j\tau_j
		=\sum_{j\notin J}
		\left(\frac{|z_j|^2}{4}e^{2\tau_j}-a_j\tau_j\right).
	\]
	最后一个等式使用 \(\sum_j\langle x,\nu_j\rangle\tau_j
	=\langle x,\beta\tau\rangle=0\)。每个括号中的一元函数在两端都趋于正无穷，且有有限下界。投影
	\(\mathfrak k\to\mathbb R^{J^c}\) 单射：若只有 \(J\) 位置可能非零，则 \(\beta\tau=0\) 与这些法向量的独立性迫使 \(\tau=0\)。因此 \(\mathcal F_z\) 在 \(\mathfrak k\) 上固有。其 Hessian 为
	\[
		\mathrm{D}^2\mathcal F_z(\tau)(\eta,\eta)
		=\sum_{j\notin J}|z_j|^2e^{2\tau_j}\eta_j^2>0
		\qquad(\eta\ne0).
	\]
	所以它有唯一极小点。临界方程是
	\(i^*(\mu_{\mathbb C}(e^\tau z)+\lambda)=0\)，恰好说明 \(e^\tau z\in Z_P\)。严格凸性给出交点的 \(K\) 轨道唯一性。

	Hessian 可逆还使极小点依赖 \(z\) 光滑，第一卷定理11.3.5的参数隐函数定理给出这一结论。在每张仿射图中选上面非活动坐标为 \(1\) 的代表，再作这个唯一正伸缩，便得到商映射的光滑逆。故所列对应为微分同胚；特别地，这些图的整体空间具有 \(M_P\) 的 Hausdorff 拓扑。

	最后在 \(z\in Z_P\) 记 \(V_z=\mathrm{T}_z(Kz)\)，并以标准 Euclid 度量取
	\[
		H_z=(V_z\oplus JV_z)^\perp.
	\]
	矩映射方程给出 \(\mathrm{T}_zZ_P=(JV_z)^\perp\)，故
	\(\mathrm{T}_zZ_P=V_z\oplus H_z\)。\(H_z\) 对标准 \(J\) 不变，且在商切空间上投影为同构；复商的微分则消去 \(V_z\oplus JV_z\)，所以其复结构正是 \(H_z\) 上的 \(J\)。标准形式在 \(H_z\) 上与此复结构相容，约化拉回等式使 \(\omega_P\) 也相容。复结构由全纯单项式过渡得到，因而可积；\(\omega_P\) 已闭，故为 Kähler 形式。
\end{proof}

这里的仿射图及粘合只依赖法扇。下一小节将用射影闭包证明这个复流形是代数簇。它具有开稠密的 \(T_{\mathbb C}\) 轨道，且乘法作用正则延拓，属于\term[daishuhuanmiancu]{代数环面簇}[algebraic toric variety]；本节的光滑模型也是正规簇。

\ProseParagraphBreak

第 \(j\) 个面对应子流形
\[
	D_j=\{[z]:z_j=0\}\subset M_P.
\]
在包含该面的顶点图中，它是一个坐标超平面；在其他图中则为空。因此它是光滑复余维一子簇，并等于 \(\mu_P^{-1}\) 的相应面原像。其开稠密部分是一个复环面轨道，所以不可约。称 \(D_j\) 为\term[huanmianbubianchuzi]{环面不变除子}[torus-invariant divisor]。这里除子是复余维一不可约子簇的有限整数线性组合；局部由一个非零有理函数给出零极点阶数的除子称\term[Cartierchuzi]{Cartier 除子}[Cartier divisor]，局部方程之比必须为可逆正则函数。本节中的各 \(D_j\) 由坐标函数给出，因此满足这一条件。

\begin{equation}\label{g4:eq:1205:character-divisor}
	\operatorname{div}(\chi^m)
	=\sum_{j=1}^d\langle m,\nu_j\rangle D_j.
\end{equation}
这可在每个顶点图中直接验证：\(\chi^m=\prod_i u_i^{\langle m,\nu_i\rangle}\)，指数就是相应超平面上的零阶或负的极点阶。

\subsection{极化和整多面体}
\label{g4:subsec:1205:02}

法扇不记录支撑常数的实数大小。要把某一个辛形式与射影嵌入的线丛对应起来，还需要整数条件。称顶点都在 \(M_{\mathbb Z}\) 中的多面体为\term[zhengduomianti]{整多面体}[integral polytope]。对 Delzant 多面体，这等价于全部 \(\lambda_j\) 为整数：一边由面上的整数顶点求配对，另一边由顶点的整数法向基反解顶点坐标。

\ProseParagraphBreak

这里也明确线丛术语。若一个代数线丛的一组整体截面在各点不同时为零，并给出到某个射影空间的闭嵌入，就称它为\term[henchongzuxiancong]{很充足线丛}[very ample line bundle]；某个正张量幂很充足时，称它为\term[chongzuxiancong]{充足线丛}[ample line bundle]。本节的射影模型以一条指定的充足线丛作为极化数据。下面将实际写出所需的截面和嵌入。

\begin{proposition}\label{g4:prop:1205:integral-embedding}
	设 \(P\) 为整 Delzant 多面体。则其单项式图所成的空间同构于射影闭包
	\begin{equation}\label{g4:eq:1205:projective-monomials}
		X_P=\overline{\{[\chi^{-m}(t)]_{m\in P\cap M_{\mathbb Z}}:
		t\in T_{\mathbb C}\}}\ \subset
		\mathbb{CP}^{\,|P\cap M_{\mathbb Z}|-1}.
	\end{equation}
	负指数与本章正旋转矩映射取负二次型的约定相配。除子
	\(D_P=\sum_j\lambda_jD_j\) 定义一条线丛 \(L_P\)，上述坐标由其整体截面给出，且 \(L_P\) 很充足。
\end{proposition}
\begin{proof}
	固定顶点 \(v\)，仍用其外法向对偶基 \(m_1,\ldots,m_n\)。每个 \(m\in P\cap M_{\mathbb Z}\) 满足
	\[
		v-m=\sum_i a_i m_i,\qquad
		a_i=\lambda_i-\langle m,\nu_i\rangle\in\mathbb Z_{\ge0}.
	\]
	从 \(v\) 沿第 \(i\) 条边到另一端点，位移为 \(-\ell_i m_i\)。两个端点都整，且 \(m_i\) 原始，所以 \(\ell_i\) 为正整数；特别地 \(v-m_i\in P\cap M_{\mathbb Z}\)。在射影图 \(Z_v\ne0\) 中，
	\[
		\frac{Z_m}{Z_v}=\chi^{v-m}
		=\prod_i u_i^{a_i},
		\qquad
		\frac{Z_{v-m_i}}{Z_v}=u_i.
	\]
	所以环面像在这张图中的闭包，恰为所有这些单项式的图像，且因坐标 \(u_i\) 本身都出现，它同构于 \(\mathbb C^n\)。像为闭图，逆映射是选取所列 \(n\) 个比值，因而双向正则。

	这些顶点图覆盖整个射影闭包。事实上，每个格点 \(m\in P\) 可写成顶点的有理凸组合；清去分母后，有
	\(Nm=\sum_v a_vv\)、\(a_v\in\mathbb Z_{\ge0}\)、\(\sum_v a_v=N\)。
	因此在环面像及其闭包上恒有齐次关系
	\[
		Z_m^N=\prod_v Z_v^{a_v}.
	\]
	若所有顶点坐标均为零，上式使全部坐标都为零，与射影点的定义矛盾。由此，每个点至少落在一张顶点图内。图间过渡仍是字符的 Laurent 单项式，故与前面构造的空间完全一致。

	下面落实线丛。由\eqref{g4:eq:1205:character-divisor}，\(D_P\) 在顶点图中由有理函数 \(\chi^v\) 给出，因为活动面的阶数为 \(\langle v,\nu_i\rangle=\lambda_i\)。在重叠上，这些函数之比可逆，因而可用局部标架 \(e_v\) 及过渡式
	\[
		e_w=\chi^{v-w}e_v
	\]
	定义代数线丛 \(L_P\)。对每个 \(m\in P\cap M_{\mathbb Z}\)，令截面 \(s_m\) 在 \(v\) 图中的系数为 \(\chi^{v-m}\)。刚才的非负指数保证它正则，过渡式保证各图相容。且 \(s_v\) 在 \(v\) 图中系数恒为 \(1\)，故这些截面没有共同零点。它们给出的映射正是\eqref{g4:eq:1205:projective-monomials}，而已证明的局部闭图及覆盖说明它是射影闭嵌入。这证明很充足性。
\end{proof}

即使原 \(P\) 的支撑常数为实数，其法扇仍来自某个整 Delzant 多面体。为此，把有限个 \(\lambda_j\) 作足够小的有理扰动：已有顶点的非活动不等式严格成立，有限个其他独立面交点若不属于 \(P\)，至少违反一个严格不等式；所以小扰动保持全部顶点面关联。再乘一个共同正整数，就得到具有同一法扇的整多面体。因此上述单项式图总能实现为一个光滑射影簇。原来的实数支撑常数决定其上的具体 Kähler 类。

\begin{proposition}\label{g4:prop:1205:polarization-class}
	对整 Delzant 多面体 \(P\)，在前述识别下有
	\begin{equation}\label{g4:eq:1205:polarization-class}
		c_1(L_P)_{\mathbb R}=\left[\frac{\omega_P}{2\pi}\right].
	\end{equation}
	更一般地，对任意实 Delzant 多面体，\([\omega_P/(2\pi)]\) 是整数上同调类的实像，当且仅当 \(P\) 可经平移成为整多面体。
\end{proposition}
\begin{proof}
	先设所有 \(\lambda_j\) 为整数。它们定义 \(K_{\mathbb C}\) 的字符
	\(\chi^\lambda(k)=\prod_jk_j^{\lambda_j}\)。
	在 \(\mathcal U\times\mathbb C\) 上取作用
	\(k\cdot(z,a)=(kz,\chi^\lambda(k)a)\)。其商是一条线丛。多项式
	\[
		\prod_j z_j^{\,\lambda_j-\langle m,\nu_j\rangle}
	\]
	在 \(K_{\mathbb C}\) 下按字符 \(\chi^\lambda\) 变换，故给出商线丛截面。在顶点规范化中，其系数恰为前面的 \(\chi^{v-m}\)，所以商线丛正是 \(L_P\)。

	限制到 \(Z_P\)，用 \(|a|^2\) 赋予商线丛 Hermitian 度量。在环境空间取
	\[
		\alpha_0=\frac12\sum_j(x_j\,\mathrm{d}y_j-y_j\,\mathrm{d}x_j),
		\qquad \mathrm{d}\alpha_0=\omega_0,
		\qquad \nabla_0=\mathrm{d}-i\alpha_0.
	\]
	对 \(\xi\in\mathfrak k\)，水平集方程说明
	\(\alpha_0(\xi_{\mathbb C^d})=\sum_j|z_j|^2\xi_j/2
	=\langle\lambda,\xi\rangle\)。
	纤维相位的变化与连接形式的这一项相抵消，所以连接沿 \(K\) 下降。也可直接取局部截面 \(z=z(x)\)：局部连接形式是 \(-i z^*\alpha_0\)；改成 \(z'=e^{it(x)}z\) 时，标架乘以 \(e^{-i\langle\lambda,t(x)\rangle}\)，连接形式也恰增加
	\(-i\,\mathrm{d}\langle\lambda,t(x)\rangle\)。这核对了下降和符号。下降连接的曲率为 \(-i\omega_P\)，第二卷命题22.3.1的一阶陈类曲率归一化给出\eqref{g4:eq:1205:polarization-class}。

	若 \(P\) 平移 \(a\)，则支撑常数变成
	\(\lambda+\beta^*a\)，其在 \(\mathfrak k\) 上的限制不变。因此水平集、约化形式都不变，矩映射仅加上 \(a\)。这证明平移后整时的充分性。

	反过来，设 \([\omega_P/(2\pi)]\) 整。考虑任一边
	\(E=[v,v+\ell u]\)，其中 \(u\in M_{\mathbb Z}\) 为原始内向边向量，\(\ell>0\)。由顶点图与\cref{g4:prop:1203:orbit-fibres}，其原像是一条二维球面：开边上为一个圆族，两个端点的圆分别缩成一点。取与 \(u\) 配对为 \(-1\) 的剩余原始圆生成元，令其周期 \(2\pi\) 的参数为 \(\theta\)。对开边参数 \(s\in(0,\ell)\)，矩映射方程给出 \(\iota_{\partial_\theta}\omega_P=-\mathrm{d}s\)，故限制面积形式为
	\(\mathrm{d}s\wedge\mathrm{d}\theta\)，选此辛定向后
	\[
		\int_{\mu_P^{-1}(E)}\frac{\omega_P}{2\pi}=\ell.
	\]
	故每条边的格长度为整数。顶点边图连通：取顶点值互异的线性函数，从非最大顶点总有一条边使其增加，因为这些内向边生成该顶点的可行方向锥；沿这样的边有限步就达到唯一最大顶点。因此所有顶点与一个固定顶点之差都是格向量。把这个固定顶点平移到零，便得到整多面体。
\end{proof}

射影性与指定辛类的整性由此分开。相同法扇的实多面体可以给出同一个射影簇上的不同实 Kähler 类；只有相应平移整条件成立时，这个类才由一条极化线丛的一阶陈类给出。对射影直线，\(\int\omega_{\mathrm{FS}}=\pi\)，所以 \(\mathcal O(1)\) 对应的是 \(2\omega_{\mathrm{FS}}\) 的类，而不是 \(\omega_{\mathrm{FS}}\) 的类。

\subsection{射影空间、乘积及吹起模型}
\label{g4:subsec:1205:03}

对单纯形
\[
	P_s=\{x_i\ge0,\ \sum_{i=1}^n x_i\le s\},\qquad s>0,
\]
外法向量为 \(-e_1,\ldots,-e_n,(1,\ldots,1)\)，支撑常数为 \(0,\ldots,0,s\)。核环面是对角圆，水平方程为 \(\sum_{j=0}^n|z_j|^2=2s\)，所以
\begin{equation}\label{g4:eq:1205:simplex-projective}
	(M_{P_s},\omega_{P_s})\cong
	(\mathbb{CP}^n,2s\,\omega_{\mathrm{FS}}).
\end{equation}
其复商是 \((\mathbb C^{n+1}\setminus\{0\})/\mathbb C^*\)，正是射影空间的比例坐标。整数 \(s\) 时，格点对应总次数不超过 \(s\) 的仿射单项式，齐次化后为全部 \(s\) 次齐次单项式，因此极化为 \(\mathcal O(s)\)。

\ProseParagraphBreak

若 \(P=P_1\times P_2\)，法向映射、核群和约束方程按两个因子分开。故
\[
	M_P=M_{P_1}\times M_{P_2},\qquad
	\omega_P=\operatorname{pr}_1^*\omega_{P_1}
	+\operatorname{pr}_2^*\omega_{P_2}.
\]
复商图也按乘积粘合。特别地，盒子 \(\prod_i[0,a_i]\) 对应
\((\mathbb{CP}^1)^n\) 及形式 \(\sum_i2a_i\operatorname{pr}_i^*\omega_{\mathrm{FS}}\)。
在整情形，格点集合为笛卡尔积，射影坐标是两组单项式坐标的两两乘积；这就是第五卷命题1.3.3的 Segre 嵌入，\footnote{Corrado Segre（1863--1924），意大利数学家，研究高维射影几何与代数曲面。}极化线丛为两因子极化的外张量积。

\ProseParagraphBreak

再考虑在顶点 \(v\) 截角。设活动外法向量为 \(\nu_1,\ldots,\nu_n\)，对偶基为 \(m_i\)。加上新不等式
\begin{equation}\label{g4:eq:1205:corner-cut}
	\left\langle x,\sum_{i=1}^n\nu_i\right\rangle
	\le\sum_{i=1}^n\lambda_i-\varepsilon.
\end{equation}
当 \(\varepsilon>0\) 足够小时，它只截去顶点附近的单纯形，产生新顶点 \(v-\varepsilon m_i\)。新外法向量 \(\sum_i\nu_i\) 原始；在每个新顶点，它替换原基中的一个向量，行列式仍为 \(\pm1\)，所以截角后仍为 Delzant 多面体。

\begin{proposition}\label{g4:prop:1205:corner-blowup}
	设 \(n\ge2\)。在上述足够小截角中，\(M_{P_\varepsilon}\) 是 \(M_P\) 在顶点固定点处的辛吹起。局部的复模型为
	\[
		\operatorname{Bl}_0\mathbb C^n
		=\{(Z,[\ell])\in\mathbb C^n\times\mathbb{CP}^{n-1}:Z\in\ell\}.
	\]
	被替换的点成为\term[liwaichuzi]{例外除子}[exceptional divisor]
	\(\mathbb{CP}^{n-1}\)，其辛形式为 \(2\varepsilon\omega_{\mathrm{FS}}\)。
\end{proposition}
\begin{proof}
	先将\cref{g4:prop:1203:vertex-chart}的标准辛坐标两两合成复数 \(z_i\)，并记
	\(I_i=|z_i|^2/2\)。截角条件正是 \(\sum_iI_i\ge\varepsilon\)。在
	\(\mathbb C^n_z\times\mathbb C_\eta\) 上取标准形式，并让圆作用为
	\[
		t\cdot(z,\eta)=(e^{-it}z,e^{it}\eta),\qquad
		q(z,\eta)=\frac12|z|^2-\frac12|\eta|^2.
	\]
	在 \(q=\varepsilon>0\) 上，\(z\ne0\)，作用自由。若 \(\eta\ne0\)，可唯一把它转为正实数，截面为
	\[
		z\longmapsto\bigl(z,\sqrt{|z|^2-2\varepsilon}\bigr),
		\qquad |z|^2>2\varepsilon.
	\]
	额外复坐标的虚部恒为零，因此约化形式在这部分严格等于原来的标准形式。若 \(\eta=0\)，剩下
	\(|z|^2=2\varepsilon\) 除以标量圆，得到
	\((\mathbb{CP}^{n-1},2\varepsilon\omega_{\mathrm{FS}})\)。
	这就是把一个球的内部移去、把边界的相位圆压成点的辛吹起局部模型。

	给这个标准局部模型配上通常的复结构，也可直接描述其代数商。复化作用为
	\(\rho\cdot(z,\eta)=(\rho^{-1}z,\rho\eta)\)，\(z\ne0\)。
	映射
	\[
		(z,\eta)\longmapsto(Z=\eta z,[z])
	\]
	沿复轨道恒定，且诱导到 \(\operatorname{Bl}_0\mathbb C^n\) 的双射。
	在 \(z_i\ne0\) 的图中把 \(z_i\) 规范成 \(1\)，坐标为
	\(a=\eta z_i\)、\(t_j=z_j/z_i\)，并且
	\(Z_i=a,\ Z_j=at_j\)。这些公式给出双向正则的仿射图。对任一这样的复轨道，正实伸缩下
	\[
		q(e^{-\tau}z,e^\tau\eta)
		=\tfrac12e^{-2\tau}|z|^2-\tfrac12e^{2\tau}|\eta|^2
	\]
	严格递减，并唯一经过 \(\varepsilon>0\)，故这个复商也与辛约化的光滑空间一致。以上辛截面中的 \(z\) 与代数吹降中的 \(Z=\eta z\) 分别给出了外部辛识别和全纯吹降映射。

	最后核对它确为截角多面体的构造。新法向映射是
	\(\beta'(e_j)=\nu_j\)、\(\beta'(e_{d+1})=\sum_{i=1}^n\nu_i\)。
	其核群含原来的 \(K\)，剩余圆由整数向量
	\((-\sum_{i=1}^n e_i,e_{d+1})\) 生成。先按 \(K\) 约化，得到
	\(M_P\times\mathbb C\)；剩余圆的零水平方程为
	\[
		\sum_{i=1}^n h_i(\mu_P)-\frac12|\eta|^2-\varepsilon=0.
	\]
	它在顶点图内正是 \(q=\varepsilon\)。当 \(\eta\ne0\) 时作用自由；当 \(\eta=0\) 时，足够小的 \(\varepsilon\) 使点留在上述顶点图内，仍自由。因此分阶段约化定理\cref{g4:thm:1003:reduction-stages}把这个局部替换与 \(M_{P_\varepsilon}\) 识别。外部部分由刚才的截面与未截去区域辛同构。

	在法扇中，截角把顶点最大锥分成由新射线
	\(\sum_i\nu_i\) 与其余 \(n-1\) 条旧射线生成的 \(n\) 个锥。它们的单项式坐标变换为
	\(a_k=a_it_k,\ t_j'=t_j/t_k,\ t_i'=1/t_k\)，恰是上面吹起图之间的变换。因此代数侧的局部替换也是同一个点吹起。
\end{proof}

二维时，\cref{g4:fig:1203:polytopes}中的新边对应例外射影直线，其面积为 \(2\pi\varepsilon\)。若原多面体已整，则保留原格并取整数 \(\varepsilon\) 才能直接得到整截角多面体；实数截角仍给出辛形式，只是所选类未必整。顶点截角与吹起的关系也见 Cannas da Silva 的命题6.9；这里的面积常数按本书已经固定的 \(2\pi\) 周期和 Fubini--Study 归一化计算。\cite[第6.4节，第172页]{Cannas2006Symplectic}

\subsection{一般非光滑多面体的新现象}
\label{g4:subsec:1205:04}

简单、有理与光滑格条件分别控制不同现象。若顶点仍恰有 \(n\) 个独立有理法向量，但它们生成的子格
\(N'\subset N\) 指数大于 \(1\)，则在 \(N'\) 中它们是一组基，而在原格中不是。相应仿射图成为有限群商
\[
	\mathbb C^n/\Gamma,\qquad
	\Gamma=\operatorname{Hom}(M'_{\mathbb Z}/M_{\mathbb Z},\mathbb C^*),
	\qquad M'_{\mathbb Z}=\operatorname{Hom}(N',\mathbb Z).
\]
确实，\(\mathbb C^n\) 的单项式指数属于对偶格 \(M'_{\mathbb Z}\)，有限群不变的单项式恰好是指数落在 \(M_{\mathbb Z}\) 中的那些。因此不变环就是原格所规定的半群环。辛侧的无穷小稳定子仍为零，但可能留下有限稳定子，对应\secref{g4:sec:1004}的局部有限商。

\begin{example}\label{g4:exm:1205:finite-quotient-cone}
	二维锥由 \(\nu_1=(1,0)\)、\(\nu_2=(1,2)\) 生成，行列式为 \(2\)。在法向量生成的子格中取标准仿射坐标 \(u,v\)，剩余二阶群同时把它们变为 \((-u,-v)\)。不变环由
	\[
		A=u^2,\qquad B=uv,\qquad C=v^2
	\]
	生成，并满足 \(AC=B^2\)。每个不变单项式总次数为偶数：两指数同为偶数时由 \(A,C\) 生成，同为奇数时再乘一个 \(B\)，这也证明没有其他生成元；用关系消去 \(B^2\) 后的单项式在 \(u,v\) 中互异，故没有其他关系。于是局部商为
	\(\{AC=B^2\}\subset\mathbb C^3\)。由第五卷定理5.3.5及例5.3.11，原点的切空间维数为 \(3\)，而局部维数为 \(2\)，所以原点奇异。整数格指数在这里直接变成有限商奇点。
\end{example}

若多面体不简单，一个顶点可能落在多于 \(n\) 个面上。假设有 \(r>n\) 个活动面，它们的法向量张成 \(\mathfrak t\)，则核群与活动坐标环面的交具有维数 \(r-n>0\)。因此约化水平在该点不正则，已不能由局部自由作用的有限群商解释；约束本身参与奇异结构。

\begin{example}\label{g4:exm:1205:nonsimple-cone}
	在 \(N=\mathbb Z^3\) 中取四条射线
	\[
		(0,0,1),\quad(1,0,1),\quad(0,1,1),\quad(1,1,1).
	\]
	它们生成一个非单纯锥，可作为一个四棱锥顶点的法锥。其对偶整半群由
	\[
		(1,0,0),\quad(0,1,0),\quad(-1,0,1),\quad(0,-1,1)
	\]
	生成。验证方法是：对偶条件为
	\(c\ge0,\ a+c\ge0,\ b+c\ge0,\ a+b+c\ge0\)，等价于
	\(c\ge\max(-a,0)+\max(-b,0)\)。取非负整数 \(C,D\)，使
	\(C+D=c,\ C\ge\max(-a,0),\ D\ge\max(-b,0)\)，再取
	\(A=a+C,\ B=b+D\)，便得到所需的非负整数分解。

	相应四个单项式记为 \(U,V,W,Z\)，唯一关系为 \(UW=VZ\)。任一两组同指数分解的差都是 \((1,-1,1,-1)\) 的整数倍，因而由这个二项式关系连接；故仿射模型为
	\[
		\{UW-VZ=0\}\subset\mathbb C^4.
	\]
	它在原点的局部维数为 \(3\)、切空间维数为 \(4\)，是一个奇异三维二次锥。这里四条法向射线只张成三维空间，正维稳定子与这一非简单顶点相伴出现。
\end{example}

另外，若面的法向方向本身不有理，就不能用整数映射 \(\beta\) 将逐坐标圆群映到固定的 \(T\)；相应实核未必给出闭子环面。若法向量有理却故意使用非原始倍数，则会保留额外的有限稳定子，这些整数标记也不包含在无标记的多面体形状中。因而一般非光滑情形需要同时记录格、稳定子和分层数据；光滑紧辛环面流形的 Delzant 条件正好使这些局部商成为普通光滑仿射图。
% 纲要条目排为编号小节标题，并进入目录与书签；说明段仍保留于纲要注释中。

% !TeX root = ../../Book4.tex
% 纲要标题：### 12.6*　等变局部化与约化体积
% 纲要位置：几何学计划纲要.md，第 3963 行。
% 正文状态：初稿；尚未正式审核。
% 写作范围、先修与边界：先修为 IV.10 的正则辛约化、IV.12.1 的固定点权和局部标准形，以及 II.附录 C 的等变上同调；本节保持选读，不反向成为凸性或 Delzant 分类的前提。固定紧连通辛流形上的 Hamilton 环面作用，以实或有理系数并对非零权局部化；完整计算孤立固定点的球面和复射影空间模型，证明圆作用自由正则水平集上约化类及体积的变化。一般固定子流形的局部化和一般分段多项式结论按精确版本外引，局部自由、轨形与奇异约化不混同于自由商。

% 纲要内容开始
% 
% > 写作范围与外部输入：先修为 IV.10 的正则辛约化、IV.12.1 的固定点权和局部标准形，以及 II.附录 C 的等变上同调；本节保持选读，不反向成为凸性或 Delzant 分类的前提。固定紧连通辛流形上的 Hamilton 环面作用，以实或有理系数并对非零权局部化；完整计算孤立固定点的球面和复射影空间模型，证明圆作用自由正则水平集上约化类及体积的变化。一般固定子流形的局部化和一般分段多项式结论按精确版本外引，局部自由、轨形与奇异约化不混同于自由商。
% 
% * Cartan 模型与等变积分的归一化
% * 固定点权、等变 Euler 类与局部化公式
% * 正则值上的约化形式与上同调类变化
% * Duistermaat–Heckman 测度的分段多项式性
% * 球面与复射影空间的局部化及体积计算
% 
% ---
% 
% 纲要内容结束

\OptionalSection{等变局部化与约化体积}{g4:sec:1206}

一个 Hamilton 环面作用同时提供两种积分方法。把积分推到矩映射的像上，可以用约化空间的体积作密度；把积分限制到固定点，则会出现切表示的权。两种方法在球面和射影空间中给出同一个答案，但其中的角周期、权与面积单位必须一致。本节恢复第二卷附录 C 所需的等变上同调约定，证明孤立固定点的积分公式及自由圆作用正则层的体积变化，并计算这些模型。一般固定分支公式与一般分段多项式定理在相应位置明确引用。

\subsection{Cartan 模型与等变积分的归一化}
\label{g4:subsec:1206:01}
设 \(T=\mathbb R^k/(2\pi\mathbb Z)^k\) 光滑作用于紧连通无边界辛流形 \((M^{2n},\omega)\)。记第 \(j\) 个角方向的基本场为 \(V_j\)，其流的参数周期为 \(2\pi\)。一个整数权 \(\alpha=(\alpha_1,\ldots,\alpha_k)\in\mathbb Z^k\) 所对应的特征标为
\[
	\chi_\alpha(\theta)=e^{\mathrm{i}\langle\alpha,\theta\rangle}.
\]
若用周期为 \(1\) 的生成元作 Lie 代数基，则基本场变成 \(2\pi V_j\)。第二卷附录 C 的整数权变量 \(u_j\) 正是这一基的对偶变量：其上同调次数为 \(2\)，拓扑上代表权 \(1\) 关联线的一阶陈类。

\ProseParagraphBreak

因此本节的\term[cartanmoxing]{Cartan 模型}[Cartan model]写为
\begin{equation}\label{g4:eq:1206:cartan-normalization}
	\Omega_T(M)=\bigl(\mathbb R[u_1,\ldots,u_k]\otimes\Omega(M)\bigr)^T,
	\qquad
	\mathrm{d}_T=\mathrm{d}-2\pi\sum_{j=1}^k u_j\iota_{V_j}.
\end{equation}
这里的 Cartan 指 Henri Cartan。\footnote{Henri Paul Cartan（1904--2008），法国数学家，研究多复变、代数拓扑与同调代数。}多项式次数按两倍计入总次数，形式次数按通常方式计入。第二卷定义 C.1.6、命题 C.1.7采用同一周期为 \(1\) 的变量约定；特别地，在不变形式上
\[
	\mathrm{d}_T^2=-2\pi\sum_j u_j\mathcal L_{V_j}=0.
\]
第二卷定理 C.1.10的紧群等变 de Rham 比较把这个复形的上同调识别为 Borel 等变上同调 \(H_T^*(M;\mathbb R)\)。\footnote{Armand Borel（1923--2003），瑞士数学家，研究代数拓扑与线性代数群；他与测度论中 Borel 集的 Félix Édouard Justin Émile Borel 是不同人物。}比较定理是本节的外部理论输入之一，其紧连通版本见 Meinrenken 的作者稿第6节定理6.1；这里 \(T\) 确为紧连通群。\cite[作者稿第6节，定理6.1]{Meinrenken2006EquivariantCartan}
% 跨卷实读：g2:def:cartan-complex(C.1.6)、g2:prop:cartan-is-complex(C.1.7)、g2:thm:equivariant-de-rham(C.1.10)及所列外引范围；AppendixC.tex第275—445行。

\begin{definition}\label{g4:def:1206:equivariant-symplectic-form}
	设上述作用为 Hamilton 作用，\(\mu=(\mu_1,\ldots,\mu_k)\) 为矩映射，满足 \(\mathrm{d}\mu_j=\iota_{V_j}\omega\)。定义归一化的\term[dengbianxinxingshi]{等变辛形式}[equivariant symplectic form]为
	\begin{equation}\label{g4:eq:1206:equivariant-omega}
		\Omega_T=\frac{\omega}{2\pi}+\sum_j\mu_j u_j.
	\end{equation}
	它是总次数为 \(2\) 的等变闭形式，因为
	\[
		\mathrm{d}_T\Omega_T
		=\sum_j u_j\bigl(\mathrm{d}\mu_j-\iota_{V_j}\omega\bigr)=0.
	\]
	令 \(u=0\) 得到普通类 \([\omega/(2\pi)]\)。这个实上同调类是否为整类是另一项条件，定义本身不要求它整。
\end{definition}

等变积分逐个多项式系数取普通积分，并且只积分形式次数为 \(2n\) 的部分：
\[
	\int_M\sum_\nu u^\nu\eta_\nu
	=\sum_\nu u^\nu\int_M(\eta_\nu)_{[2n]}.
\]
它使总次数降低 \(2n\)。若被积形式是 \(\mathrm{d}_T\)-恰当的，最高形式次数的贡献只来自普通外微分，因为插入运算需要一个不存在的 \(2n+1\) 次形式；Stokes 公式遂说明其积分为零。因此积分下降到等变上同调。

\ProseParagraphBreak

在计算中还会用到指数 \(e^{\Omega_T}\)。可以先按总次数展开，逐项使用积分公式；由于 \(M\) 紧，也可以在给定实参数 \(u\) 后把它解释为通常的指数积分。所得到的表达式是
\begin{equation}\label{g4:eq:1206:normalized-laplace}
	\int_M e^{\Omega_T}
	=\int_M e^{\langle\mu,u\rangle}
	\frac{\omega^n}{(2\pi)^n n!}.
\end{equation}
归一化因子 \((2\pi)^{-n}\) 与式\eqref{g4:eq:1206:cartan-normalization}中的 \(2\pi\) 配套，使后面的 Euler 类直接由整数权给出。

\begin{remark}\label{g4:rem:1206:curvature-sign}
	对主右圆丛 \(P\to B\)，若实连接形式 \(\alpha\) 在周期为 \(2\pi\) 的正向生成场 \(X\) 上满足 \(\alpha(X)=1\)，写 \(\mathrm{d}\alpha=\pi^*\beta\)。则第二卷式(C.10)—(C.11)的连接计算在此变成
	\[
		u\longmapsto-\frac{\beta}{2\pi},
		\qquad
		c_1(P\times_{S^1}\mathbb C)_{\mathbb R}
		=\left[-\frac{\beta}{2\pi}\right],
	\]
	其中关联线采用权 \(1\) 表示，等价关系为 \((p,z)\sim(p\cdot g,g^{-1}z)\)。理由是周期为 \(1\) 的连接为 \(\alpha/(2\pi)\)，故第二卷的 \(-\mathrm{d}a\) 恰是这里的 \(-\mathrm{d}\alpha/(2\pi)\)。这也固定了约化类变化公式中的符号。
\end{remark}

\subsection{固定点权、等变 Euler 类与局部化公式}
\label{g4:subsec:1206:02}
在固定点 \(p\) 处，选取与 \(\omega_p\) 相容且 \(T\) 不变的复结构。切表示分解成复一维权空间；若 \(p\) 孤立，所有权 \(\alpha_{p,1},\ldots,\alpha_{p,n}\) 均非零。\cref{g4:thm:1201:fixed-weight-model}把这些权实现为实际坐标中的作用
\[
	\theta\cdot(z_1,\ldots,z_n)
	=\bigl(e^{\mathrm{i}\langle\alpha_{p,1},\theta\rangle}z_1,\ldots,
	e^{\mathrm{i}\langle\alpha_{p,n},\theta\rangle}z_n\bigr),
	\qquad \omega=\sum_j\mathrm{d}x_j\wedge\mathrm{d}y_j.
\]
按第十章的正向基本场约定，局部矩映射是
\begin{equation}\label{g4:eq:1206:weight-moment}
	\mu(z)=\mu(p)-\frac12\sum_j\alpha_{p,j}|z_j|^2.
\end{equation}
权的正负由实际特征标决定；矩映射的二次项带着另一个已固定的负号。

\ProseParagraphBreak

固定点切表示的\term[dengbianeulerlei]{等变 Euler 类}[equivariant Euler class]为
\begin{equation}\label{g4:eq:1206:euler-weights}
	e_T(\mathrm{T}_pM)=\prod_{j=1}^n\langle\alpha_{p,j},u\rangle.
\end{equation}
相容复结构给出的定向与辛定向一致，所以这里不用另添定向符号。下面的局部构造同时验证这个权乘积，并给出积分公式需要的归一化。

\begin{lemma}\label{g4:lem:1206:equivariant-thom-point}
	在权为 \(\alpha\ne0\) 的复直线 \(\mathbb C\) 上，可以构造紧支等变闭形式 \(\tau_\alpha\)，使
	\[
		\int_{\mathbb C}\tau_\alpha=1,\qquad
		\tau_\alpha|_0=\langle\alpha,u\rangle.
	\]
	在 \(n\) 条权直线的直和上，存在相应的乘积形式，其积分为 \(1\)，原点限制为式\eqref{g4:eq:1206:euler-weights}。
\end{lemma}
\begin{proof}
	在 \(\mathbb C\setminus\{0\}\) 上写 \(z=re^{\mathrm{i}\vartheta}\)，令 \(\gamma=\mathrm{d}\vartheta/(2\pi)\)。各基本场在这一平面上为 \(\alpha_j\partial_\vartheta\)，故
	\[
		\mathrm{d}_T\gamma=-\langle\alpha,u\rangle.
	\]
	选光滑径向函数 \(\chi\)，在原点附近为 \(1\)，在一个圆盘之外为 \(0\)，定义
	\[
		\tau_\alpha
		=\langle\alpha,u\rangle\chi-\mathrm{d}\chi\wedge\gamma.
	\]
	第二项的支集避开原点，所以它延拓为全平面的光滑形式；第一项在原点附近为常多项式。用分次乘积法则，
	\[
		\mathrm{d}_T\tau_\alpha
		=\langle\alpha,u\rangle\,\mathrm{d}\chi
		+\mathrm{d}\chi\wedge\mathrm{d}_T\gamma=0.
	\]
	按复定向积分最高形式项，得到
	\[
		\int_{\mathbb C}\tau_\alpha
		=-\int_0^\infty\chi'(r)\,\mathrm{d}r
		\int_0^{2\pi}\frac{\mathrm{d}\vartheta}{2\pi}=1.
	\]
	原点限制则是 \(\langle\alpha,u\rangle\)。在直和中将各条直线上的形式拉回相乘，等变闭性由乘积法则保证；支集位于有界多圆盘中，积分由逐次积分等于 \(1\)，原点限制是各权的乘积。由等变 Thom 类沿零截面的限制定义，\footnote{René Thom（1923--2002），法国数学家，发展了配边理论，并研究奇点与突变理论。}这正是切表示的等变 Euler 类。
\end{proof}

允许非零权作分母后，记
\[
	R=\mathbb R[u_1,\ldots,u_k],\qquad
	S=\left\{\text{有限个非零整数权线性多项式的乘积}\right\}.
\]
第二卷定理 C.1.12给出局部化后的限制同构
\begin{equation}\label{g4:eq:1206:localized-restriction}
	S^{-1}H_T^*(M;\mathbb R)
	\ \xrightarrow{\ \cong\ }\ 
	S^{-1}H_T^*(M^T;\mathbb R).
\end{equation}
该定理的有限 \(T\)-胞腔情形已在第二卷证明；这里使用其紧光滑流形版本，原始证明见 Atiyah--Bott 第3节定理3.5。\cite[第8页，定理3.5]{AtiyahBott1984MomentMap}

\begin{theorem}[孤立固定点的等变积分]\label{g4:thm:1206:isolated-localization}
	设紧环面 \(T\) Hamilton 作用于紧无边界辛流形 \((M^{2n},\omega)\)，固定点集有限。若 \(\eta\in\Omega_T(M)\) 等变闭，则在局部化系数环中
	\begin{equation}\label{g4:eq:1206:isolated-localization}
		\int_M\eta
		=\sum_{p\in M^T}
		\frac{\eta|_p}{\prod_{j=1}^n\langle\alpha_{p,j},u\rangle}.
	\end{equation}
	特别地，在所有分母非零的实参数处，
	\begin{equation}\label{g4:eq:1206:exponential-localization}
		\int_M e^{\langle\mu,u\rangle}\frac{\omega^n}{(2\pi)^n n!}
		=\sum_{p\in M^T}
		\frac{e^{\langle\mu(p),u\rangle}}
		{\prod_{j=1}^n\langle\alpha_{p,j},u\rangle}.
	\end{equation}
\end{theorem}
\begin{proof}
	在每个固定点选取互不相交的等变坐标邻域。把上一引理的乘积形式支集取得足够小，并在邻域之外延拓为零，得到 \(M\) 上的等变闭形式 \(\tau_p\)。它在 \(p\) 的限制是 \(e_T(\mathrm{T}_pM)\)，在其他固定点限制为零，且 \(\int_M\tau_p=1\)。

	在局部化系数下，两个类
	\[
		[\eta],\qquad
		\sum_{p\in M^T}\frac{\eta|_p}{e_T(\mathrm{T}_pM)}[\tau_p]
	\]
	限制到每个固定点都相同。式\eqref{g4:eq:1206:localized-restriction}的单射性使它们相等。对两边积分，并使用 \(\int_M\tau_p=1\)，得到第一式。

	将第一式用于 \(\Omega_T^\ell/\ell!\)，再对 \(\ell\ge0\) 求和。左边只有普通形式次数为 \(2n\) 的项参与积分，恰为式\eqref{g4:eq:1206:normalized-laplace}；右边在固定点处只有零次形式，求和成为指数。对固定且避开分母零点的实参数，指数级数在紧 \(M\) 上一致收敛，右边只有有限个固定点，故求和与积分均可交换。
\end{proof}

当固定点不是孤立点时，法方向的 Euler 类仍可逆。设 \(F\) 为一个固定分支，其相容复法丛分解成权子丛
\(\nu_F=\bigoplus_{\alpha\ne0}E_\alpha\)，其中 \(r_\alpha=\operatorname{rank}_{\mathbb C}E_\alpha\)。在 \(F\) 的普通上同调与多项式环的张量积中，
\begin{equation}\label{g4:eq:1206:normal-euler}
	e_T(\nu_F)=
	\prod_{\alpha\ne0}
	\sum_{j=0}^{r_\alpha}c_j(E_\alpha)\langle\alpha,u\rangle^{r_\alpha-j},
	\qquad c_0(E_\alpha)=1.
\end{equation}
这个表达式由第二卷定理21.3.6和命题21.3.10给出：复丛的 Euler 类等于最高陈类，而在 \(F\times BT\) 上，权子丛对应于 \(E_\alpha\) 与权线丛的张量积；后者的第一陈类为 \(\langle\alpha,u\rangle\)。在线丛扭转公式中取最高陈类，即得到各个因子的表达式。零次普通上同调部分是非零权的乘积。除去这一可逆部分后，其余项含正次普通上同调类，因维数有限而幂零；所以有限几何级数就能写出逆元。
% 跨卷实读：g2:thm:top-class-euler(21.3.6)、g2:prop:chern-operations(21.3.10)及其完整证明，后者取线丛扭转最高陈类。

\begin{theorem}[等变积分局部化]\label{g4:thm:1206:fixed-component-localization}
	设紧环面 \(T\) 光滑作用于紧定向无边界流形 \(M\)。对每个固定连通分支 \(F\) 选择不变法丛定向，并由 \(\mathrm{T}M|_F=\mathrm{T}F\oplus\nu_F\) 规定 \(F\) 的定向。则对任意等变上同调类 \(\eta\)，在非零权局部化后有
	\begin{equation}\label{g4:eq:1206:fixed-component-localization}
		\int_M\eta=\sum_{F\subset M^T}\int_F\frac{\eta|_F}{e_T(\nu_F)}.
	\end{equation}
	在 Hamilton 辛流形情形，可使用固定分支及相容复法丛的辛定向。
\end{theorem}

这是 Atiyah--Bott 第3节式(3.8)的积分公式，其证明在局部化同构之外使用一般等变 Thom 推前及自交公式；第二卷式(C.12)已按同一范围引用。原文第3节采用复系数，本式的实系数版本由实类复化后同一等式及复化的单射性得到。本节上面的显式构造证明了孤立固定点这一特殊情形；一般法丛的推前构造取自该文第2—3节。\cite[第5、8--9页，式(2.19)、(3.7)--(3.8)]{AtiyahBott1984MomentMap}

\subsection{正则值上的约化形式与上同调类变化}
\label{g4:subsec:1206:03}
固定矩映射的一个值以后，正则辛约化给出一个较低维的辛流形。现在让这个值变化。流形可以沿正则区间保持同一微分同胚类型，而辛形式的上同调类会记录主圆丛的曲率。

\begin{theorem}\label{g4:thm:1206:reduced-class-variation}
	设 \(S^1=\mathbb R/(2\pi\mathbb Z)\) Hamilton 作用于紧无边界辛流形 \((M^{2n},\omega)\)，矩映射为 \(\mu:M\to\mathbb R\)。设 \(a\) 为正则值，且作用在非空水平集 \(P=\mu^{-1}(a)\) 上自由。记 \(B=P/S^1\)、约化形式为 \(\omega_a\)。则在充分小的 \(|t|<\varepsilon\) 内，各水平集上的作用仍自由，约化空间可光滑识别为同一个 \(B\)。把 \(P\) 的右作用定义为 \(p\cdot e^{\mathrm{i}\theta}=e^{\mathrm{i}\theta}\cdot p\)，并令
	\[
		e=c_1(P\times_{S^1}\mathbb C)_{\mathbb R}\in H^2(B;\mathbb R).
	\]
	在上述识别下有
	\begin{equation}\label{g4:eq:1206:dh-class}
		\left[\frac{\omega_{a+t}}{2\pi}\right]
		=\left[\frac{\omega_a}{2\pi}\right]+t e.
	\end{equation}
\end{theorem}
\begin{proof}
	先构造保持圆作用的水平集识别。取圆作用不变的黎曼度量；其存在由第一卷定理17.3.4的平均度量构造保证。因 \(P\) 紧且 \(\mathrm{d}\mu\) 在其上非零，紧性使我们能选 \(\varepsilon>0\)，使 \(\mu^{-1}([a-2\varepsilon,a+2\varepsilon])\) 内没有临界点。在此带状区域定义不变向量场
	\[
		Y=\frac{\nabla\mu}{|\nabla\mu|^2},
		\qquad \mathrm{d}\mu(Y)=1.
	\]
	其积分曲线满足 \(\mu(\Phi_t^Yp)=\mu(p)+t\)。当起点在 \(P\)、\(|t|\le\varepsilon\) 时，轨道始终留在上述紧区域内。第一卷定理13.2.2的有限端点逃逸判据因而保证所需的流全部存在。逆流把邻近水平集送回 \(P\)，得到等变微分同胚
	\[
		\Phi:P\times(-\varepsilon,\varepsilon)
		\longrightarrow\mu^{-1}((a-\varepsilon,a+\varepsilon)),
		\qquad \mu\Phi=a+t.
	\]
	等变微分同胚保持稳定子，所以邻近层也自由。第十章的正则辛约化定理遂在所有这些层适用。

	记拉回的辛形式为 \(\Omega=\Phi^*\omega\)，圆作用生成场仍记为 \(X\)。于是
	\[
		\iota_X\Omega=\mathrm{d}t.
	\]
	在 \(P\) 上取不变连接形式 \(\alpha\)，满足 \(\alpha(X)=1\)；例如对不变度量可取 \(\alpha=g(X,\cdot)/g(X,X)\)。将其拉回到乘积上。由不变性和 Cartan 公式，\(\mathrm{d}\alpha\) 水平且不变，因此有唯一闭形式 \(\beta\in\Omega^2(B)\) 满足 \(\mathrm{d}\alpha=\pi^*\beta\)。\cref{g4:rem:1206:curvature-sign}给出 \(e=[-\beta/(2\pi)]\)。

	考虑
	\[
		\Omega+\mathrm{d}(t\alpha).
	\]
	它闭且不变，而
	\(\iota_X\mathrm{d}(t\alpha)=-\mathrm{d}t\)，所以也水平。\cref{g4:prop:1002:basic-descent}所证明的基本形式下降使它来自 \(B\times(-\varepsilon,\varepsilon)\) 上的一个闭二形式 \(\eta\)。写
	\[
		\eta=\eta_t+\mathrm{d}t\wedge\gamma_t,
		\qquad
		\eta_t=\omega_{a+t}+t\beta.
	\]
	闭性给出 \(\mathrm{d}_B\eta_t=0\)、\(\partial_t\eta_t=\mathrm{d}_B\gamma_t\)。对 \(t\) 积分可知
	\[
		\eta_t-\eta_0
		=\mathrm{d}_B\int_0^t\gamma_s\,\mathrm{d}s,
	\]
	故 \([\omega_{a+t}]+t[\beta]=[\omega_a]\)。除以 \(2\pi\) 并代入 \(e\) 即得结论。
\end{proof}

这个证明直接控制上同调类，不要求所选流把每个约化形式本身变成线性函数。更换连接使 \(\beta\) 只改变一个恰当形式，更换由不变横向流给出的识别则产生同伦的识别；式\eqref{g4:eq:1206:dh-class}因此记录了正则区间上同一族约化空间的类变化。此处的曲率计算与 Cannas da Silva 的命题6.14、定理6.15相应；其圆群在连接计算中采用周期 \(1\)，本节已按周期 \(2\pi\) 写出全部因子。\cite[第6.6节，第179--180页]{Cannas2006Symplectic}

\begin{corollary}\label{g4:cor:1206:reduced-volume-polynomial}
	在上一条定理的条件下，令 \(v=[\omega_a/(2\pi)]\)。约化空间的归一化体积为
	\begin{equation}\label{g4:eq:1206:reduced-volume-polynomial}
		V(a+t)\coloneqq
		\int_B\frac{\omega_{a+t}^{n-1}}{(2\pi)^{n-1}(n-1)!}
		=\sum_{j=0}^{n-1}
		\frac{t^j}{j!(n-1-j)!}\int_B v^{n-1-j}e^j.
	\end{equation}
	所以它在这个自由正则区间内为次数至多 \(n-1\) 的多项式。
\end{corollary}
\begin{proof}
	积分只依赖闭形式的上同调类：将一个因子改变为恰当形式，与其余闭因子的乘积仍恰当，在紧无边界 \(B\) 上积分为零。故可用 \((v+te)^{n-1}\) 计算，再用二项式展开得到公式。
\end{proof}

\subsection{Duistermaat–Heckman 测度的分段多项式性}
\label{g4:subsec:1206:04}
矩映射把 \(M\) 上的辛体积送到底空间上的测度。对 Borel 集 \(A\subset\mathbb R^k\)，定义\term[duistermaatheckmancesu]{Duistermaat--Heckman 测度}[Duistermaat--Heckman measure]，\footnote{Johannes Jisse Duistermaat（1942--2010），荷兰数学家，研究微局部分析与辛几何。Gerrit Jacobus Heckman（1953年生），荷兰数学家，研究表示论、辛几何与超几何函数。}记为
\begin{equation}\label{g4:eq:1206:dh-measure}
	m_{\mathrm{DH}}(A)=\int_{\mu^{-1}(A)}\frac{\omega^n}{n!}.
\end{equation}
也记归一化测度 \(\nu_{\mathrm{DH}}=(2\pi)^{-n}m_{\mathrm{DH}}\)。对连续函数 \(f\)，推前的含义是
\[
	\int_{\mathbb R^k}f(x)\,\mathrm{d}\nu_{\mathrm{DH}}(x)
	=\int_M f(\mu)\frac{\omega^n}{(2\pi)^n n!}.
\]
因此式\eqref{g4:eq:1206:exponential-localization}正给出这个归一化测度的指数积分。

\begin{proposition}\label{g4:prop:1206:free-circle-density}
	在\cref{g4:thm:1206:reduced-class-variation}的自由正则区间上，归一化测度的密度正是约化空间的归一化体积：
	\begin{equation}\label{g4:eq:1206:free-circle-density}
		\mathrm{d}\nu_{\mathrm{DH}}(a+t)=V(a+t)\,\mathrm{d}t.
	\end{equation}
	等价地，未归一化测度的密度为
	\(2\pi\int_B\omega_{a+t}^{n-1}/(n-1)!\)。
\end{proposition}
\begin{proof}
	沿用上一证明的乘积坐标与连接。由 \(\eta=\Omega+\mathrm{d}(t\alpha)\) 的分解，
	\[
		\Omega
		=\pi^*\omega_{a+t}+\alpha\wedge\mathrm{d}t
		+\mathrm{d}t\wedge\pi^*\gamma_t.
	\]
	底流形 \(B\) 的维数为 \(2n-2\)。展开最高次幂时，纯水平 \(2n\) 次项为零，含两次 \(\mathrm{d}t\) 的项也为零；含 \(\gamma_t\) 而不含 \(\alpha\) 的项在 \(B\) 方向超过维数。因此
	\[
		\frac{\Omega^n}{n!}
		=\pi^*\!\left(\frac{\omega_{a+t}^{n-1}}{(n-1)!}\right)
		\wedge\alpha\wedge\mathrm{d}t.
	\]
	正向圆轨道上 \(\int\alpha=2\pi\)。在主丛局部平凡化中逐次积分，再用底空间单位分解相加，得到
	\[
		\int_{\mu^{-1}(a+E)}\frac{\omega^n}{n!}
		=2\pi\int_E
		\left(\int_B\frac{\omega_{a+t}^{n-1}}{(n-1)!}\right)\mathrm{d}t
	\]
	对区间内的每个 Borel 集 \(E\) 成立。除以 \((2\pi)^n\) 即得所述密度。
\end{proof}

上述密度计算与类变化结合，已经完整证明自由圆作用正则层附近的多项式性。在一个全部水平集都自由的连通正则区间内，各个局部多项式在重叠开区间相等，故是同一个多项式的限制。

\begin{theorem}[Duistermaat--Heckman 定理]\label{g4:thm:1206:dh-general}
	设紧环面 \(T^k\) 有效地 Hamilton 作用于紧连通无边界辛流形 \((M^{2n},\omega)\)，矩映射为 \(\mu:M\to\mathbb R^k\)。则 \(m_{\mathrm{DH}}\) 关于 \(\mathbb R^k\) 上的 Lebesgue 测度绝对连续，并且存在一个密度，其限制到每个正则值连通区域上为多项式。对临界值集合上密度代表的取值不作规定。
\end{theorem}

此处引用 Cannas da Silva 第6.6节定理6.17；该节在第178页假定矩映射固有，紧 \(M\) 自动满足这一条件。定理在第181页陈述，第182—183页给出圆群证明概述并把一般环面情形指向原始论文及进一步文献。本节已补齐自由圆群情形的证明，一般环面及非自由正则层的结论按上述定理外引。\cite[第6.6节，定理6.17，第178、181--183页]{Cannas2006Symplectic}

\ProseParagraphBreak

有效性使矩映射所用的目标空间与实际作用方向一致。如果存在正维子环面处处平凡作用，则其矩映射分量为常数，测度集中在一个真仿射子空间中，不能相对于原来整个 \(\mathbb R^k\) 的 Lebesgue 测度保持绝对连续。另一方面，即使作用有效，正则水平集上的稳定子也可能是非平凡有限群。此时轨道商按第十章的讨论应作为轨形处理；式\eqref{g4:eq:1206:free-circle-density}中按普通自由主丛积分得到的常数，不能直接移植到这样的层。

\subsection{球面与复射影空间的局部化及体积计算}
\label{g4:subsec:1206:05}
先计算面积为 \(2\pi\lambda\)、\(\lambda>0\) 的射影直线。由\cref{g4:thm:1003:projective-reduction}，其形式为
\[
	\omega_\lambda=2\lambda\,\omega_{\mathrm{FS}},
	\qquad \int_{\mathbb CP^1}\omega_\lambda=2\pi\lambda.
\]
令圆群按
\[
	e^{\mathrm{i}\theta}\cdot[z_0:z_1]=[z_0:e^{-\mathrm{i}\theta}z_1]
\]
作用。其矩映射与像为
\[
	x=\lambda\frac{|z_1|^2}{|z_0|^2+|z_1|^2},\qquad 0\le x\le\lambda.
\]
在两极之间取 \(\vartheta=-\arg(z_1/z_0)\)，则正向圆作用使 \(\vartheta\) 增加，且 \(\omega_\lambda=\mathrm{d}\vartheta\wedge\mathrm{d}x\)。这同时给出 \(\iota_{\partial_\vartheta}\omega_\lambda=\mathrm{d}x\) 及归一化推前测度 \(\mathrm{d}x\)。

\begin{example}\label{g4:exm:1206:sphere-localization}
	固定点 \(p_0=[1:0]\) 的局部坐标 \(z_1/z_0\) 按 \(e^{-\mathrm{i}\theta}\) 变化，所以其权为 \(-1\)；\(p_1=[0:1]\) 的局部坐标 \(z_0/z_1\) 按 \(e^{\mathrm{i}\theta}\) 变化，权为 \(1\)。两点的矩映射值分别为 \(0,\lambda\)。因此
	\[
		\int_{\mathbb CP^1}e^{xu}\frac{\omega_\lambda}{2\pi}
		=\frac1{-u}+\frac{e^{\lambda u}}u
		=\frac{e^{\lambda u}-1}{u}.
	\]
	直接积分也给出 \(\int_0^\lambda e^{xu}\,\mathrm{d}x\)，与局部化完全一致。右端在 \(u=0\) 的奇性可去，极限为 \(\lambda\)，恢复面积 \(2\pi\lambda\)。展开幂级数还得到全部矩：
	\[
		\int_{\mathbb CP^1}x^r\frac{\omega_\lambda}{2\pi}
		=\frac{\lambda^{r+1}}{r+1},\qquad r=0,1,2,\ldots.
	\]
	每个内点的约化空间为一点，归一化零维体积为 \(1\)，所以自由圆群密度公式也给出同一个常密度。
\end{example}

在 \(\mathbb CP^n\) 上继续取 \(\omega_\lambda=2\lambda\omega_{\mathrm{FS}}\)，其中 \(\omega_{\mathrm{FS}}\) 严格采用第十章式\eqref{g4:eq:1003:projective-form}的归一化。令
\[
	(\theta_1,\ldots,\theta_n)\cdot[z_0:z_1:\cdots:z_n]
	=[z_0:e^{-\mathrm{i}\theta_1}z_1:\cdots:e^{-\mathrm{i}\theta_n}z_n].
\]
由球面约化中的坐标函数 \(|z_j|^2/2\) 下降，矩映射为
\begin{equation}\label{g4:eq:1206:projective-moment}
	x_j=\lambda\frac{|z_j|^2}{\sum_{\ell=0}^n|z_\ell|^2},
	\qquad
	\Delta_\lambda=\left\{x_j\ge0,\ \sum_{j=1}^n x_j\le\lambda\right\}.
\end{equation}
每个满足这些不等式的点都可由 \(|z_j|^2/( \sum|z_\ell|^2)=x_j/\lambda\) 及剩余比例 \(1-\sum_jx_j/\lambda\) 实现，故像恰为该单纯形。

\begin{proposition}\label{g4:prop:1206:projective-localization}
	设 \(n\ge1\)、\(\lambda>0\)，并采用上述作用和形式。记 \(u_0=0\)。在 \(u_0,u_1,\ldots,u_n\) 两两不同的实参数处，
	\begin{equation}\label{g4:eq:1206:simplex-laplace}
		\int_{\Delta_\lambda}e^{\sum_{j=1}^n x_ju_j}\,
		\mathrm{d}x_1\cdots\mathrm{d}x_n
		=\sum_{i=0}^n
		\frac{e^{\lambda u_i}}{\prod_{j\ne i}(u_i-u_j)}.
	\end{equation}
	归一化测度 \(\nu_{\mathrm{DH}}\) 是 \(\Delta_\lambda\) 上的 Lebesgue 测度，且
	\begin{equation}\label{g4:eq:1206:projective-volume}
		\int_{\mathbb CP^n}\frac{\omega_\lambda^n}{n!}
		=(2\pi)^n\frac{\lambda^n}{n!}.
	\end{equation}
\end{proposition}
\begin{proof}
	固定点是 \(n+1\) 个坐标点 \(p_i=[0:\cdots:1:\cdots:0]\)。令 \(\varepsilon_j\) 为 \(\mathbb Z^n\) 的标准基，并约定 \(\varepsilon_0=0\)。在 \(p_i\) 处用 \(z_j/z_i\)、\(j\ne i\)，作局部坐标；分子和分母的字符相除，得到实际权
	\[
		\alpha_{i,j}=\varepsilon_i-\varepsilon_j,\qquad j\ne i.
	\]
	因此
	\[
		e_T(\mathrm{T}_{p_i}\mathbb CP^n)=\prod_{j\ne i}(u_i-u_j),
		\qquad \langle\mu(p_i),u\rangle=\lambda u_i.
	\]
	代入\cref{g4:thm:1206:isolated-localization}即得到指数积分的右端。

	再从约化形式直接计算推前测度。在所有齐次坐标均非零的稠密开集，选球面代表
	\[
		z_0=\sqrt{2\left(\lambda-\sum_jx_j\right)},\qquad
		z_j=\sqrt{2x_j}\,e^{-\mathrm{i}\vartheta_j}\quad(j\ge1).
	\]
	这里 \(x\in\operatorname{int}\Delta_\lambda\)，\(\vartheta_j\) 周期为 \(2\pi\)。第十章计算约化形式时使用的环境原始形式
	\(\lambda_0=\frac12\sum_\ell(x_\ell^{\mathrm{amb}}\,\mathrm{d}y_\ell^{\mathrm{amb}}-y_\ell^{\mathrm{amb}}\,\mathrm{d}x_\ell^{\mathrm{amb}})\)
	在这个截面上拉回为 \(-\sum_jx_j\,\mathrm{d}\vartheta_j\)。因 \(\mathrm{d}\lambda_0=\omega_0\)，得到
	\[
		\omega_\lambda=\sum_j\mathrm{d}\vartheta_j\wedge\mathrm{d}x_j.
	\]
	积分角变量后，每个角圆贡献 \(2\pi\)，恰好消去归一化因子。被排除的坐标超平面是有限个实余维二子流形，对光滑体积测度为零。因此 \(\nu_{\mathrm{DH}}\) 在整个像上就是单纯形 Lebesgue 测度，证明式\eqref{g4:eq:1206:simplex-laplace}的左端及其与右端的相等。

	单纯形的 Euclid 体积为 \(\lambda^n/n!\)：逐次对最后一个坐标积分，维数归纳给出
	\[
		\operatorname{Vol}(\Delta_\lambda)
		=\int_0^\lambda\frac{(\lambda-s)^{n-1}}{(n-1)!}\,\mathrm{d}s
		=\frac{\lambda^n}{n!}.
	\]
	乘回 \((2\pi)^n\) 得到辛体积。
\end{proof}

局部化右端各项单独在参数相等时有极点，整体却代表一个光滑积分。最低阶的消去也可以纯代数地核验：对 \(r<n\)，多项式 \(X^r\) 的插值公式中 \(X^n\) 的系数为零；对 \(r=n\) 则为 \(1\)。因此
\[
	\sum_{i=0}^n\frac{u_i^r}{\prod_{j\ne i}(u_i-u_j)}
	=\begin{cases}0,&0\le r<n,\\1,&r=n.\end{cases}
\]
这里插值公式是
\(X^r=\sum_i u_i^r\prod_{j\ne i}(X-u_j)/(u_i-u_j)\)，两边次数至多 \(n\) 且在 \(n+1\) 个不同点取值相同，所以相等。将其用于 \(\Omega_T^n\)，便独立读出 \(\int_{\mathbb CP^n}(\omega_\lambda/(2\pi))^n=\lambda^n\)。固定点权、作用角体积和单纯形体积由此给出了三种相容的计算。

% 纲要条目排为编号小节标题，并进入目录与书签；说明段仍保留于纲要注释中。


\begin{chaptersummary}
环面作用在固定点附近由整数权表示，在一般轨道附近由完整稳定子、轨道余切方向及辛切片共同表示。矩映射分量的临界集是相应子环面的固定集；非零权使法向 Hessian 非退化，正负法丛均为复丛。跨越临界值时附着的是负法丛的盘丛，指标及余指标的偶性由此保证分量水平集连通。

\ProseParagraphBreak

完整矩纤维的连通性还需要同时控制所有分量。将每条纤维的连通分支压成一点，局部标准形给出凸锥图；紧商上的最短路在这些图中只能映为直线段。因此商映射到矩像单射，矩纤维连通且矩像凸。环面秩的归纳在极值分支中找到共同固定点，最终得到有限固定点像的凸包公式。

\ProseParagraphBreak

有效半维作用使每个稳定子都为连通环面，切片权构成整数基，矩纤维恰是一条轨道。矩像因而满足 Delzant 条件。反过来，对每个面引入一个复坐标，令其模长平方等于相应支撑函数的两倍，再按法向量关系作辛约化，便恢复该多面体。顶点的整数基条件排除有限稳定子，给出普通光滑流形。

\ProseParagraphBreak

同一个多面体给出相同的局部边界模型。保持矩映射的局部等变辛自同构可由底空间上的光滑 Hamilton 函数表示；这些函数的核是常整数仿射函数。凸多面体上相应的 Čech 误差可以消去，随后用单位分解粘成整体同构，完成分类的唯一性。忘掉矩映射常数对应多面体平移，重标环面对应整数线性变换，尺度则保留为辛几何数据。

\ProseParagraphBreak

法扇决定单项式坐标与不变除子，支撑常数决定 Kähler 类；平移后为整多面体，恰对应辛类 \([\omega/(2\pi)]\) 为整类。选读部分以同一个整数权约定计算等变积分：孤立固定点贡献为限制类除以切权乘积，自由圆群正则约化的辛类随矩值线性变化。球面与射影空间的固定点计算和多面体体积给出一致的答案。
\end{chaptersummary}

\BookExercises

\begin{exercise}\label{g4:ex:12:mixed-weights-fixed-line}
	让 \(S^1=\mathbb R/(2\pi\mathbb Z)\) 在标准辛空间 \(\mathbb C^3\) 上作用为
	\[
		\theta\cdot(z_1,z_2,z_3)
		=(e^{2\mathrm i\theta}z_1,e^{-3\mathrm i\theta}z_2,z_3).
	\]
	取在原点为零的矩函数。求固定集、全部可能的稳定子、临界子流形的指标与余指标，并显式描述正、负、零矩纤维的连通性。说明固定子流形的切向零特征值为什么不违背 Morse–Bott 非退化性。
\end{exercise}
\begin{solution}
	正旋转权给出
	\[
		\mu=-|z_1|^2+\frac32|z_2|^2.
	\]
	固定集为 \(C=\{z_1=z_2=0\}\cong\mathbb C\)。若 \(z_1\neq0,z_2=0\)，稳定子为二次单位根群 \(\mu_2\)；若 \(z_1=0,z_2\neq0\)，稳定子为 \(\mu_3\)；若两者都非零，由 \(\gcd(2,3)=1\) 知稳定子平凡；两者都为零时稳定子为整个圆群。\(z_3\) 不影响这些条件。

	法向实坐标中的 Hessian 为 \(\operatorname{diag}(-2,-2,3,3)\)，所以指标与余指标都等于 \(2\)。核恰为 \(\mathrm TC\)，法向 Hessian 非退化，这正是 Morse–Bott 条件；它并不要求整个环境 Hessian 非退化。

	当 \(a>0\) 时，由 \(\mu=a\) 得
	\[
		z_2=\sqrt{\tfrac23(a+|z_1|^2)}\,e^{\mathrm i\phi},
	\]
	所以 \((z_1,e^{\mathrm i\phi},z_3)\) 给出 \(\mu^{-1}(a)\cong\mathbb C^2\times S^1\)。当 \(a<0\) 时，可用
	\[
		z_1=\sqrt{\tfrac32|z_2|^2-a}\,e^{\mathrm i\phi}
	\]
	得到同样的微分同胚类型。当 \(a=0\) 时，写 \(z_2=\rho e^{\mathrm i\phi_2}\)、\(z_1=\sqrt{3/2}\rho e^{\mathrm i\phi_1}\)，\(\rho\geq0\)。在 \(\rho=0\) 处全部相位合并，因此零纤维同胚于 \(\operatorname{Cone}(S^1\times S^1)\times\mathbb C\)，其中锥参数为 \([0,\infty)\)。三种纤维都连通。这个例子有非平凡有限稳定子，也有非孤立临界集，仍具有权模型所描述的偶法向指标。
\end{solution}

\begin{exercise}\label{g4:ex:12:finite-stabilizer-model}
	固定整数 \(m\geq2\)。在 \(M=S^1_\varphi\times\mathbb R_\eta\times\mathbb C_w\) 上置 \(w=x+\mathrm iy\)、\(I=|w|^2/2\)，并定义
	\[
		\Omega=\frac1m\mathrm d\varphi\wedge\mathrm d(\eta+I)
		+\mathrm dx\wedge\mathrm dy.
	\]
	证明 \(\Omega\) 是辛形式，作用
	\[
		\theta\cdot(\varphi,\eta,w)
		=(\varphi+m\theta,\eta,e^{\mathrm i\theta}w)
	\]
	有效且以 \(\mu=\eta\) 为矩函数。求稳定子，并用映射
	\[
		[t,\eta,z]\longmapsto(t^m,\eta,tz)
	\]
	把这个空间与 \(S^1\times_{\mu_m}(\mathbb R\times\mathbb C)\) 的全稳定子模型作等变辛识别。这里 \(\mu_m\) 在 \(\mathbb C\) 上按通常乘法作用。
\end{exercise}
\begin{solution}
	形式闭，且
	\[
		\Omega^2=\frac2m\mathrm d\varphi\wedge\mathrm d\eta
		\wedge\mathrm dx\wedge\mathrm dy
	\]
	处处非零，故它非退化。基本场为 \(X=m\partial_\varphi-y\partial_x+x\partial_y\)，并满足 \(XI=0\)。直接收缩得
	\[
		\iota_X\Omega=\mathrm d(\eta+I)-\mathrm dI=\mathrm d\eta.
	\]
	所以矩函数为 \(\eta\)。在 \(w=0\) 时，固定条件为 \(m\theta\in2\pi\mathbb Z\)，稳定子正是 \(\mu_m\)；在 \(w\neq0\) 时还需 \(e^{\mathrm i\theta}=1\)，稳定子平凡。后一个条件也说明整个作用有效。

	商的等价关系为 \((t,\eta,z)\sim(th^{-1},\eta,hz)\)，其中 \(h^m=1\)。因此 \(t^m\) 与 \(tz\) 都与代表元无关。每个 \((e^{\mathrm i\varphi},\eta,w)\) 可选择一个 \(m\) 次根 \(t\)，再取 \(z=t^{-1}w\)；不同根恰好给出同一等价类。局部选取光滑根便得到光滑逆，所以这是等变微分同胚。

	为核对形式，在商前的覆盖上写 \(t=e^{\mathrm i\theta}\)。有 \(\varphi=m\theta\)、\(w=e^{\mathrm i\theta}z\) 及 \(I_w=I_z\)。由极坐标计算，或在零点外计算后光滑延拓，
	\[
		\mathrm dx_w\wedge\mathrm dy_w
		=\mathrm dx_z\wedge\mathrm dy_z
		+\mathrm dI_z\wedge\mathrm d\theta.
	\]
	故所给映射把 \(\Omega\) 拉回为
	\[
		\mathrm d\theta\wedge\mathrm d\eta
		+\mathrm d\theta\wedge\mathrm dI_z
		+\mathrm dx_z\wedge\mathrm dy_z
		+\mathrm dI_z\wedge\mathrm d\theta
		=\mathrm d\theta\wedge\mathrm d\eta
		+\mathrm dx_z\wedge\mathrm dy_z,
	\]
	正是 \(\mathfrak h=0\)、完整稳定子 \(H=\mu_m\) 时的模型形式。把 \(H\) 仅换成 \(H^0=\{1\}\) 会遗漏零截面的有限稳定子。
\end{solution}

\begin{exercise}\label{g4:ex:12:noncompact-annular-image}
	令 \(B=\{b\in\mathbb R^2:1<\|b\|<2\}\)，在 \(M=\mathbb T^2\times B\) 上取
	\[
		\omega=\mathrm d\theta_1\wedge\mathrm db_1
		+\mathrm d\theta_2\wedge\mathrm db_2,
	\]
	并让 \(\mathbb T^2\) 平移角坐标。证明 \(\mu(\theta,b)=b\) 是矩映射，其每条非空纤维连通，而像不凸。再求任意非零线性投影 \(\langle\mu,\xi\rangle\) 的像，并找一个不连通的分量水平集。
\end{exercise}
\begin{solution}
	\(\omega\) 闭且非退化，\(\iota_{\partial_{\theta_j}}\omega=\mathrm db_j\)，所以 \(\mu=b\)。每条纤维为 \(\mathbb T^2\times\{b\}\)，紧且连通；但像为开圆环 \(B\)，例如两个相反点之间的线段经过不在 \(B\) 中的原点，所以像不凸。

	若 \(\xi\neq0\)，Cauchy 不等式给出投影像包含于 \((-2\|\xi\|,2\|\xi\|)\)。反过来，对区间内任一数 \(c\)，取半径 \(r\) 满足 \(\max\{1,|c|/\|\xi\|\}<r<2\)，在半径 \(r\) 的圆上可找到与 \(\xi\) 配对为 \(c\) 的向量。因此投影像恰为这个区间。然而分量 \(b_1\) 的零水平对应
	\[
		\{b_1=0,\ 1<|b_2|<2\}
		=\bigl(\{0\}\times(-2,-1)\bigr)
		\sqcup\bigl(\{0\}\times(1,2)\bigr),
	\]
	乘上 \(\mathbb T^2\) 后仍有两个连通分支。这里总流形不紧，故不满足凸性定理或紧 Morse–Bott 水平连通定理的假设。所有一维投影的像都是区间，也确实没有迫使二维矩像凸。
\end{solution}

\begin{exercise}\label{g4:ex:12:boundary-phase}
	在穿孔辛平面 \((\mathbb C^*,\mathrm{d}x\wedge\mathrm{d}y)\) 上令圆群作正旋转，矩函数为 \(\mu(z)=-|z|^2/2\)。证明
	\[
		F(z)=e^{\mathrm{i}/|z|^2}z
	\]
	是保持矩函数的圆等变辛同胚。它能否连续、光滑地延至原点？将 \(F\) 写成某个 \(f\circ\mu\) 的时间 \(1\) Hamilton 流，并解释生成函数在矩像端点处的行为。
\end{exercise}
\begin{solution}
	令 \(z=re^{\mathrm{i}\theta}\)，则 \(F\) 为 \((r,\theta)\mapsto(r,\theta+r^{-2})\)。它与圆旋转交换，保持 \(r\) 及 \(\mu\)，逆映射把角增量改为 \(-r^{-2}\)。又因为
	\[
		F^*(r\,\mathrm{d}r\wedge\mathrm{d}\theta)
		=r\,\mathrm{d}r\wedge\bigl(\mathrm{d}\theta-2r^{-3}\mathrm{d}r\bigr)
		=r\,\mathrm{d}r\wedge\mathrm{d}\theta,
	\]
	它是辛同胚。

	由 \(|F(z)|=|z|\)，令 \(F(0)=0\) 即得连续延拓，逆映射同样连续。但沿正实轴有 \(F(r)/r=e^{\mathrm{i}/r^2}\)。取 \(r_j=(2\pi j)^{-1/2}\) 与 \(s_j=((2j+1)\pi)^{-1/2}\)，该比值分别恒为 \(1\) 和 \(-1\)，所以方向导数不存在，不能微分，更不能光滑。

	对 \(b<0\) 取
	\[
		f(b)=-\frac12\log(-b),\qquad f'(b)=-\frac1{2b}.
	\]
	\(X_{f\circ\mu}=f'(\mu)X_\mu=r^{-2}\partial_\theta\)，其时间 \(1\) 流正是 \(F\)。当 \(b\uparrow0\) 时，\(f\) 及其导数发散。即使允许在连通内部给 \(f\) 加上核函数 \(c+2\pi mb\)，也不能消去这一发散。因此内部的辛同构并不自动满足\cref{g4:prop:1204:boundary-automorphisms}的边界光滑假设。
\end{solution}

\begin{exercise}\label{g4:ex:12:rectangle-rigidity}
	设 \(a,b,c,d>0\)，矩形 \(P=[0,a]\times[0,b]\)、\(Q=[0,c]\times[0,d]\) 均采用标准整数格。判定何时其辛环面流形在允许重标环面的意义下等变辛同构。计算其四维辛体积，并据此说明相同体积是否足以判定这种同构。
\end{exercise}
\begin{solution}
	由 Delzant 分类及\cref{g4:prop:1204:translation-lattice-freedom}，问题等价于是否存在整数仿射变换 \(x\mapsto Ax+t\)、\(A\in\operatorname{GL}(2,\mathbb Z)\)，把 \(P\) 送到 \(Q\)。它把顶点及相邻边送到顶点及相邻边。原始格向量在 \(A\) 下仍原始：若 \(Au=kv\)，则 \(u=kA^{-1}v\)，与原始性矛盾。目标矩形的原始边向量只有 \(\pm e_1,\pm e_2\)，所以 \(A\) 必把两条坐标轴分别送到这两条轴，允许交换及改变方向。它保留每条边的格长度，故必要条件为
	\[
		\{a,b\}=\{c,d\}
	\]
	作为带重数的无序对相等。此条件也充分：恒等或交换坐标的整数矩阵就把两个矩形相互对应，必要时再平移。

	相应流形是两个球面的乘积，面积分别为 \(2\pi a\)、\(2\pi b\)。记两因子的形式为 \(\omega_1,\omega_2\)，则 \(\omega_1^2=\omega_2^2=0\)，所以
	\[
		\int_{M_P}\frac{(\omega_1+\omega_2)^2}{2}
		=\int_{M_P}\omega_1\wedge\omega_2=(2\pi)^2ab.
	\]
	例如边长 \((1,6)\) 与 \((2,3)\) 的体积相同，却不满足无序边长对相等，故不能在允许环面重标的意义下等变辛同构。这个结论针对带环面作用的分类，所用论证没有另行判定忘掉作用后的任意辛同胚问题。
\end{solution}

\begin{exercise}\label{g4:ex:12:ineffective-same-image}
	在单位球面上用极点外的角度 \(\theta\) 和高度 \(z\in(-1,1)\) 表示面积形式 \(\mathrm{d}\theta\wedge\mathrm{d}z\)。对每个正整数 \(k\)，取
	\[
		\omega_k=\frac1{2k}\,\mathrm{d}\theta\wedge\mathrm{d}z,
		\qquad e^{\mathrm{i}t}\cdot(\theta,z)=(\theta+kt,z).
	\]
	证明各作用都有矩映射 \(\mu_k=(z+1)/2\)，矩像均为 \([0,1]\)，但不同 \(k\) 的辛流形不同构。计算作用的核，并说明该例为何不违背 Delzant 唯一性。
\end{exercise}
\begin{solution}
	给定形式是通常球面面积形式的正常数倍，故在极点处光滑且非退化。基本场为 \(k\partial_\theta\)，于是
	\[
		\iota_{k\partial_\theta}\omega_k=\frac12\mathrm{d}z
		=\mathrm{d}\mu_k.
	\]
	矩映射在两极的值为 \(0,1\)，所以每个矩像都是指定区间。作用核为 \(\{e^{2\pi\mathrm{i}j/k}:0\le j<k\}\)，当且仅当 \(k=1\) 时作用有效。

	各球面的辛面积为
	\[
		\int_{S^2}\omega_k
		=\frac1{2k}\int_0^{2\pi}\int_{-1}^1\mathrm{d}z\,\mathrm{d}\theta
		=\frac{2\pi}{k}.
	\]
	辛同胚保存面积，故不同 \(k\) 的球面连忘掉作用后的辛同胚也不存在。Delzant 唯一性要求固定环面有效作用，本例在 \(k>1\) 时缺少这一假设。改用有效商圆群时，若仍按周期 \(2\pi\) 归一化基本场，矩函数要改为 \((z+1)/(2k)\)，其区间长度变为 \(1/k\)，恰好恢复面积等于 \(2\pi\) 乘格长度的关系。
\end{solution}

\begin{exercise}\label{g4:ex:1205:irrational-box-class}
	设 \(a,b>0\)，考虑矩形 \(P=[0,a]\times[0,b]\) 的 Delzant 流形。识别其复流形，写出约化辛形式，并判定以下条件：\([\omega_P/(2\pi)]\) 整；存在 \(c>0\)，使 \([c\omega_P/(2\pi)]\) 整。特别讨论 \(a=1,b=\sqrt2\)，说明“复流形可射影嵌入”与“当前辛形式经某个正比例缩放后来自极化”之间的区别。
\end{exercise}
\begin{solution}
	矩形是两个区间的乘积，所以复流形为
	\(\mathbb{CP}^1\times\mathbb{CP}^1\)，约化形式为
	\[
		\omega_P=2a\,\operatorname{pr}_1^*\omega_{\mathrm{FS}}
		+2b\,\operatorname{pr}_2^*\omega_{\mathrm{FS}}.
	\]
	两个因子球面的面积分别为 \(2\pi a,2\pi b\)。因此整性必要求 \(a,b\in\mathbb Z\)。反之，若它们都是正整数，外张量积线丛
	\(\mathcal O(a)\boxtimes\mathcal O(b)\) 给出极化，其一阶陈类为上述辛类，所以条件也充分。

	缩放后同理要求 \(ca,cb\) 都为正整数。这蕴含 \(a/b=(ca)/(cb)\in\mathbb Q\)。反过来，若 \(a/b=r/s\)，其中 \(r,s\) 为互素正整数，取 \(c=r/a=s/b\) 即可。因此所求正比例缩放存在，当且仅当 \(a/b\) 有理。

	对 \(a=1,b=\sqrt2\)，这个比值无理，故不存在这样的 \(c\)。但复流形仍由 Segre 映射
	\[
		([z_0:z_1],[w_0:w_1])
		\longmapsto[z_0w_0:z_0w_1:z_1w_0:z_1w_1]
	\]
	嵌入 \(\mathbb{CP}^3\)，其像为一个光滑二次曲面。这个嵌入使用 \(\mathcal O(1)\boxtimes\mathcal O(1)\)，并不代表给定的无理比例辛类。
\end{solution}

\begin{exercise}\label{g4:ex:1205:trapezoid-normal-lines}
	设 \(a,b>0\)、\(k\in\mathbb Z_{\ge0}\)，令
	\[
		P_{a,b,k}=\{(x,y):0\le x\le a,\ 0\le y\le b+kx\}.
	\]
	证明它是 Delzant 多边形。按 \(x=0,x=a,y=0,y=b+kx\) 的顺序编号各面，求法向映射的整数核基、核环面作用及约化方程。再证明不变除子 \(D_1=\mu^{-1}(x=0)\)、\(D_2=\mu^{-1}(x=a)\) 都是射影直线，计算其面积，并用两个单项式图确定它们的复法线丛分别为 \(\mathcal O(-k)\)、\(\mathcal O(k)\)。
\end{exercise}
\begin{solution}
	原始外法向量及支撑常数分别为
	\[
		\nu_1=(-1,0),\quad\nu_2=(1,0),\quad
		\nu_3=(0,-1),\quad\nu_4=(-k,1),
		\qquad \lambda=(0,a,0,b).
	\]
	四个顶点为 \((0,0),(a,0),(a,b+ka),(0,b)\)。逐顶点取相邻两列，其行列式分别为 \(\pm1\)，故满足光滑格条件。法向映射为
	\[
		\beta(t_1,t_2,t_3,t_4)=(-t_1+t_2-kt_4,-t_3+t_4).
	\]
	整数核基可取 \((1,1,0,0)\)、\((0,k,1,1)\)，所以
	\[
		(s,t)\cdot(z_1,z_2,z_3,z_4)
		=(e^{is}z_1,e^{i(s+kt)}z_2,e^{it}z_3,e^{it}z_4).
	\]
	记 \(I_j=|z_j|^2/2\)，约化方程为
	\[
		I_1+I_2=a,\qquad kI_2+I_3+I_4=ka+b.
	\]
	矩坐标则为 \(x=I_1\)、\(y=I_3\)，其余两式给出
	\(I_2=a-x\)、\(I_4=b+kx-y\)。

	在 \(D_1\) 上，\(z_1=0\)、\(z_2\ne0\)，且
	\((z_3,z_4)\ne(0,0)\)。复核群作用可先把 \(z_2\) 规范成 \(1\)，剩余作用只按共同相位识别 \((z_3,z_4)\)，因此 \(D_1\cong\mathbb{CP}^1\)，齐次坐标为 \([z_3:z_4]\)。同样，\(D_2\) 上 \(z_2=0,z_1\ne0\)，也得到同一个射影坐标。两条矩边的格长度分别为 \(b\)、\(b+ka\)，故
	\[
		\int_{D_1}\omega=2\pi b,\qquad
		\int_{D_2}\omega=2\pi(b+ka).
	\]

	为计算 \(D_1\) 的法线，在 \(z_2z_3\ne0\) 的图中取
	\[
		u=z_4/z_3,\qquad
		\eta_3=z_1z_3^k/z_2;
	\]
	在 \(z_2z_4\ne0\) 的图中取
	\(v=z_3/z_4=1/u\)、\(\eta_4=z_1z_4^k/z_2\)。
	两图中 \(D_1\) 分别由 \(\eta_3=0,\eta_4=0\) 给出，且
	\[
		\eta_4=u^k\eta_3.
	\]
	这正是 \(\mathcal O(-k)\) 的纤维坐标变换：重言线
	\(\mathcal O(-1)\) 的两个代表向量 \((1,u)\)、\((v,1)\)
	使纤维系数满足 \(a_4=ua_3\)，取 \(k\) 次张量幂便得上式。因此 \(D_1\) 的法线丛为 \(\mathcal O(-k)\)。

	对 \(D_2\)，在 \(z_1z_3\ne0\) 与 \(z_1z_4\ne0\) 处分别取
	\[
		\zeta_3=\frac{z_2}{z_1z_3^k},\qquad
		\zeta_4=\frac{z_2}{z_1z_4^k}.
	\]
	其变换为 \(\zeta_4=u^{-k}\zeta_3\)，故法线丛为对偶的
	\(\mathcal O(k)\)。当 \(k=0\) 时，两条法线丛平凡，核作用也分成两个互不耦合的圆作用，恢复射影直线的乘积模型。
\end{solution}

\begin{exercise}\label{g4:ex:1205:two-corner-cuts}
	设 \(s,\varepsilon,\delta>0\)，且 \(\varepsilon+\delta<s\)。在三角形
	\[
		P_s=\{x,y\ge0,\ x+y\le s\}
	\]
	的顶点 \((0,0)\)、\((s,0)\) 分别作格长度为 \(\varepsilon,\delta\) 的截角。求新多边形的半空间表示、全部顶点和边的格长度，验证 Delzant 条件，并计算五条不变射影直线的辛面积。判定其辛类何时整。最后说明 \(\varepsilon+\delta\uparrow s\) 时哪条曲线的面积趋于零，为什么该边界值不能按同样方式保留这五条不变辛曲线。
\end{exercise}
\begin{solution}
	第一个截角增加 \(x+y\ge\varepsilon\)。在顶点 \((s,0)\)，两个活动外法向量为 \((0,-1),(1,1)\)，其和为 \((1,0)\)，所以第二个截角增加 \(x\le s-\delta\)。新多边形为
	\[
		\{x,y\ge0,\ \varepsilon\le x+y\le s,\ x\le s-\delta\}.
	\]
	按逆时针顺序，顶点是
	\[
		(\varepsilon,0),\quad(s-\delta,0),\quad
		(s-\delta,\delta),\quad(0,s),\quad(0,\varepsilon).
	\]
	全部边长都为正。每条边的原始方向分别可取
	\((1,0),(0,1),(-1,1),(0,-1),(1,-1)\)；
	故其格长度按这个顺序为
	\[
		s-\varepsilon-\delta,\qquad \delta,\qquad
		s-\delta,\qquad s-\varepsilon,\qquad \varepsilon.
	\]
	相邻原始外法向量的行列式均为 \(\pm1\)，也可由两次截角时的基底替换直接验证 Delzant 条件。

	每条边的原像为不变射影直线，其面积是格长度的 \(2\pi\) 倍。因此五个面积依次为
	\[
		2\pi(s-\varepsilon-\delta),\quad 2\pi\delta,\quad
		2\pi(s-\delta),\quad2\pi(s-\varepsilon),\quad2\pi\varepsilon.
	\]
	两个新边分别给出例外射影直线，面积为
	\(2\pi\varepsilon\)、\(2\pi\delta\)。若辛类整，这两个边长以及 \(s-\delta\) 都为整数，所以 \(s,\varepsilon,\delta\) 全为整数。反之，这三个参数都整时，所有顶点整，极化类公式给出辛类整。因此这是充要条件。

	当 \(\varepsilon+\delta\uparrow s\) 时，底边 \(y=0\) 的格长度趋于零，它的原像曲线面积趋于零。这条曲线是原三角形底边的两端分别被截去后的严格变换。在极限值处它被压成一点；一条仍作为辛子流形保留的闭曲线必须有严格正面积，所以这个极限不再属于上述保留全部不变曲线的辛约化模型。若删去届时冗余的 \(y\ge0\) 不等式，可重新讨论较少面的多边形构造，但那已经改变了被保留的不变曲线及流形模型。
\end{solution}

\begin{optionalexercise}\label{g4:ex:1206:ineffective-sphere}
	在面积为 \(2\pi\lambda\)、\(\lambda>0\) 的 \(\mathbb CP^1\) 上，令 \(S^1=\mathbb R/(2\pi\mathbb Z)\) 按
	\[
		e^{\mathrm{i}\theta}\cdot[z_0:z_1]
		=[z_0:e^{-m\mathrm{i}\theta}z_1],
		\qquad m\ge2
	\]
	作用。求矩映射、两固定点的切权及一般点的稳定子；分别用固定点公式与变量代换计算归一化 Duistermaat--Heckman 测度。说明把一般约化轨道直接当作“体积为一的点”代入自由圆群密度公式会漏掉什么因子。
\end{optionalexercise}
\begin{solution}
	若正文中权 \(-1\) 作用的矩为 \(x\in[0,\lambda]\)，新基本场是原来的 \(m\) 倍，故矩映射可取 \(h=mx\)，像为 \([0,m\lambda]\)。两固定点的切权分别为 \(-m,m\)。一般点的稳定子由 \(e^{-m\mathrm{i}\theta}=1\) 给出，恰为 \(m\) 阶单位根群；作用因此有有限核，但仍有非空正则值区间。

	归一化面积 \(\omega/(2\pi)\) 对 \(x\) 的推前是 \(\mathrm{d}x\)，故对 \(h=mx\) 的推前为
	\[
		\mathrm{d}\nu_{\mathrm{DH}}(h)
		=\frac1m\mathbf1_{[0,m\lambda]}(h)\,\mathrm{d}h.
	\]
	固定点公式独立给出
	\[
		\int e^{hu}\,\mathrm{d}\nu_{\mathrm{DH}}
		=\frac1{-mu}+\frac{e^{m\lambda u}}{mu}
		=\frac{e^{m\lambda u}-1}{mu},
	\]
	与上述密度的指数积分相同。其总质量仍为 \(\lambda\)，未归一化总面积仍为 \(2\pi\lambda\)。

	这里一条实际轨道在参数 \(\theta\) 从 \(0\) 走到 \(2\pi\) 时被覆盖了 \(m\) 次。若连接满足 \(\alpha(X)=1\)，沿一次几何轨道的积分是 \(2\pi/m\)，不是自由作用时的 \(2\pi\)。所以自由主丛密度公式中的“约化点贡献一”须在此结合稳定子阶计入 \(1/m\)。这既解释了漏掉的因子，也说明有限无效核与正维无效子群不同：前者仍允许关于直线 Lebesgue 测度的密度。
\end{solution}

\begin{optionalexercise}\label{g4:ex:1206:projective-circle-density}
	在 \((\mathbb CP^2,2\lambda\omega_{\mathrm{FS}})\) 上，取圆作用
	\[
		e^{\mathrm{i}\theta}\cdot[z_0:z_1:z_2]
		=[z_0:e^{-\mathrm{i}\theta}z_1:e^{-2\mathrm{i}\theta}z_2].
	\]
	以正文的单纯形坐标表示其矩映射 \(h=x_1+2x_2\)。求三个固定点的权对，计算归一化指数积分，并从三角形中的线性投影直接求出归一化测度的分段密度。核对其总质量等于 \(\lambda^2/2\)，并指出内点正则层上的作用是否都自由。
\end{optionalexercise}
\begin{solution}
	三个坐标点的矩映射值为 \(0,\lambda,2\lambda\)。把环面权 \(\varepsilon_i-\varepsilon_j\) 限制到方向 \((1,2)\)，得到
	\[
		p_0:\ (-1,-2),\qquad
		p_1:\ (1,-1),\qquad
		p_2:\ (2,1).
	\]
	因此
	\[
		\int_{\mathbb CP^2}e^{hu}\frac{\omega_\lambda^2}{(2\pi)^2\,2!}
		=\frac1{2u^2}-\frac{e^{\lambda u}}{u^2}
		+\frac{e^{2\lambda u}}{2u^2}.
	\]
	为直接求密度，在三角形
	\(\Delta_\lambda=\{x_1,x_2\ge0,\ x_1+x_2\le\lambda\}\)
	中改用坐标 \((h,y)=(x_1+2x_2,x_2)\)。其逆为 \(x_1=h-2y,x_2=y\)，Jacobian 行列式为 \(1\)。不等式变成
	\[
		\max(0,h-\lambda)\le y\le h/2.
	\]
	所以
	\[
		f(h)=
		\begin{cases}
			h/2,&0<h<\lambda,\\
			(2\lambda-h)/2,&\lambda<h<2\lambda,\\
			0,&h<0\text{ 或 }h>2\lambda.
		\end{cases}
	\]
	端点处任选连续延拓值即可；由推前的定义，\(\mathrm{d}\nu_{\mathrm{DH}}=f(h)\,\mathrm{d}h\)。对两段分别积分，
	\[
		\int f(h)\,\mathrm{d}h
		=\frac{\lambda^2}{4}+\frac{\lambda^2}{4}
		=\frac{\lambda^2}{2}.
	\]
	对 \(\int f(h)e^{hu}\,\mathrm{d}h\) 分段分部积分，得到前面的三项固定点表达式；其 \(u=0\) 处的可去极限也等于 \(\lambda^2/2\)。

	只要 \(z_0,z_2\ne0,z_1=0\)，该点的稳定子便是 \(\{1,-1\}\)，并且其矩映射值可以遍历整个开区间 \((0,2\lambda)\)。因此除去临界值 \(\lambda\) 后，两段正则区间上的层也都含有非自由轨道；所得密度虽完全可以用普通三角形积分计算，却不能直接用正文自由主圆丛公式把商当作无稳定子的光滑约化空间。
\end{solution}

\begin{optionalexercise}\label{g4:ex:1206:surface-reduction-slope}
	设一个四维紧 Hamilton 圆流形满足\cref{g4:thm:1206:reduced-class-variation}的条件，并设参考约化空间 \(B\) 连通。以 \(e\) 表示该定理中的关联线第一陈类，令
	\[
		A(t)=\int_B\omega_{a+t},\qquad
		m=\int_B e.
	\]
	证明 \(m\in\mathbb Z\) 且 \(A(t)=A(0)+2\pi mt\)。计算区间 \([a+s,a+t]\) 上方区域的四维辛体积，其中 \(-\varepsilon<s<t<\varepsilon\)。再证明：若 \(e=0\in H^2(B;\mathbb R)\)，则充分邻近的约化辛曲面彼此辛同胚；若 \(m\ne0\)，不同参数的约化曲面不可能辛同胚。
\end{optionalexercise}
\begin{solution}
	紧连通约化曲面 \(B\) 由辛形式定向，因而有整数基本类。\(e\) 是整数第一陈类在实上同调中的像，所以它与基本类的配对 \(m\) 为整数。对约化类变化公式积分，立即得到
	\[
		A(t)=A(0)+2\pi mt.
	\]
	其成立区间还必须满足 \(A(0)+2\pi mt>0\)，因为约化形式保持辛定向。

	未归一化的 Duistermaat--Heckman 密度在四维情形为 \(2\pi A(r)\)。因此所求辛体积为
	\[
		\int_{\mu^{-1}([a+s,a+t])}\frac{\omega^2}{2}
		=2\pi A(0)(t-s)+2\pi^2m(t^2-s^2).
	\]
	区间端点的水平集为余维一光滑子流形，对四维体积测度为零，所以闭区间或开区间给出同一积分。

	若 \(e=0\)，正文证明中 \([\beta]=0\)，可取一形式 \(\zeta\) 使 \(\beta=\mathrm{d}\zeta\)。由那里的闭形式分解，
	\[
		\partial_t\omega_{a+t}
		=\mathrm{d}\gamma_t-\beta
		=\mathrm{d}(\gamma_t-\zeta).
	\]
	这是一个光滑原始一形式族。对任意充分小的闭参数区间，\(B\) 紧，且 \(\omega_{a+t}\) 全程辛；\cref{g4:thm:0801:moser}遂给出同位，使拉回的约化形式恒等于参考形式。因此邻近约化曲面辛同胚。若 \(m\ne0\)，则不同参数的面积严格不同，而辛同胚保持辛面积，所以它们不可能辛同胚。
\end{solution}
